% GENERATED by scripts/gen_sources_block.py -- do not edit by hand.
% 153 cited sources; 152 carry captured verbatim quotes;
% sources without quotes were registered from retrieved pages whose text
% did not yield a clean quotable sentence -- cite them from the fetched page.
\noindent [6] Carnitine - Health Professional Fact Sheet - Office of Dietary. \textit{NIH Office of Dietary Supplements}. \url{https://ods.od.nih.gov/factsheets/carnitine-HealthProfessional} (accessed 2026-11-15).
\noindent\hspace*{2em}\small\textit{Evidence:} ``Healthy children and adults do not need to consume carnitine from food or supplements because the liver and kidneys synthesize sufficient amounts to meet daily needs''
\noindent\hspace*{2em}\small\textit{Evidence:} ``Dietary carnitine has a bioavailability of about 63\% to 75\%''
\noindent\hspace*{2em}\small\textit{Evidence:} ``Absorption of supplemental L-carnitine is about 14\% to 18\%, much less than that of dietary L-carnitine''
\noindent [14] Wall BT; Stephens FB; Constantin-Teodosiu D; Marimuthu K; Macdonald IA; Greenhaff PL (2011). Chronic oral ingestion of l-carnitine and carbohydrate increases muscle .... \textit{The Journal of Physiology} 589(Pt 4):963-973. \href{https://doi.org/10.1113/jphysiol.2010.201343}{doi:10.1113/jphysiol.2010.201343}. \url{https://pmc.ncbi.nlm.nih.gov/articles/PMC3060373} (accessed 2026-11-15).
\noindent\hspace*{2em}\small\textit{Evidence:} ``After 30 years of endeavour, this is the first study to show that muscle carnitine content can be increased in humans by dietary means''
\noindent\hspace*{2em}\small\textit{Evidence:} ``after 24 weeks of supplementation perceived exertion was lower in the Carnitine group when compared to baseline (14.0 \_vs.\_ 15.0, respectively; \_P\_ < 0.05) and Control at 24 weeks (14.0 \_vs.\_ 16.2, respectively; \_P\_ < 0.05).''
\noindent\hspace*{2em}\small\textit{Evidence:} ``Collectively, these metabolic effects resulted in a reduced perception of effort and increased work output during a validated exercise performance test.''
\noindent [18] Yarizadh H; Shab-Bidar S; Zamani B; Vanani AN; Baharlooi H; Djafarian K (2020). The Effect of L-Carnitine Supplementation on Exercise .... \textit{J Am Coll Nutr} 39(5):457-468. \href{https://doi.org/10.1080/07315724.2019.1661804}{doi:10.1080/07315724.2019.1661804}. \url{https://pubmed.ncbi.nlm.nih.gov/32154768} (accessed 2026-11-15).
\noindent\hspace*{2em}\small\textit{Evidence:} ``Pooled data from seven studies showed that L-Carnitine resulted in significant improvements in muscle soreness (MS) at the five follow-up time points (0, 24, 48, 72 and 96 hours (h)) compared to placebo. Also, pooled data indicated that L-Carnitine significantly reduced creatine kinase (CK), myoglobin (Mb), and lactate dehydrogenase (LDH) levels at one follow-up period (24 h). However, no effects have been observed beyond this period.''
\noindent\hspace*{2em}\small\textit{Evidence:} ``Overall, L-carnitine supplementation ameliorated muscle damage only in resistance training groups and untrained population.''
\noindent [65] Drug Interactions between levocarnitine and warfarin. \textit{Drugs.com}. \url{https://www.drugs.com/drug-interactions/levocarnitine-with-warfarin-1454-0-2311-0.html?professional=1} (accessed 2026-11-15).
\noindent\hspace*{2em}\small\textit{Evidence:} ``Limited data suggest that levocarnitine may potentiate the hypoprothrombinemic effects of warfarin and other coumarin anticoagulants.''
\noindent\hspace*{2em}\small\textit{Evidence:} ``There have been isolated case reports of INR increase and/or bleeding in patients treated with warfarin or acenocoumarol following the addition of levocarnitine therapy.''
\noindent\hspace*{2em}\small\textit{Evidence:} ``INR levels returned to normal after discontinuation of levocarnitine.''
\noindent [66] U.S. Food and Drug Administration / DailyMed (2024). Label: LEVOCARNITINE injection, solution. \textit{DailyMed, National Library of Medicine}. \url{https://www.dailymed.nlm.nih.gov/dailymed/drugInfo.cfm?setid=2685d8da-010e-40d7-914f-80546634888d} (accessed 2026-11-15).
\noindent\hspace*{2em}\small\textit{Evidence:} ``For the prevention and treatment of carnitine deficiency in patients with end stage renal disease who are undergoing dialysis.''
\noindent\hspace*{2em}\small\textit{Evidence:} ``The recommended starting dose is 10 to 20 mg/kg dry body weight as a slow 2 to 3 minute bolus injection into the venous return line after each dialysis session.''
\noindent\hspace*{2em}\small\textit{Evidence:} ``End Stage Renal Disease (ESRD) patients on maintenance hemodialysis may have low plasma carnitine concentrations and an increased ratio of acylcarnitine/carnitine because of reduced intake of meat and dairy products, reduced renal synthesis and dialytic losses.''
\noindent [69] Label: LEVOCARNITINE tablet. \textit{U.S. National Library of Medicine (DailyMed)}. \url{https://www.dailymed.nlm.nih.gov/dailymed/drugInfo.cfm?setid=0f749812-8f02-4eeb-91d9-5dd6fc9fa159} (accessed 2026-11-15).
\noindent\hspace*{2em}\small\textit{Evidence:} ``Various mild gastrointestinal complaints have been reported during the long-term administration of oral L- or D,L-carnitine; these include transient nausea and vomiting, abdominal cramps, and diarrhea.''
\noindent\hspace*{2em}\small\textit{Evidence:} ``Decreasing the dosage often diminishes or eliminates drug-related patient body odor or gastrointestinal symptoms when present.''
\noindent\hspace*{2em}\small\textit{Evidence:} ``Tolerance should be monitored very closely during the first week of administration, and after any dosage increases.''
\noindent [76] Steele CN; Baugh ME; Griffin LE; Neilson AP; Davy BM; Hulver MW; Davy KP (2021). Fasting and postprandial trimethylamine N‐oxide in .... \textit{Physiological reports} 9(16):e14970. \href{https://doi.org/10.14814/phy2.14970}{doi:10.14814/phy2.14970}. \url{https://pmc.ncbi.nlm.nih.gov/articles/PMC8371342} (accessed 2026-11-15).
\noindent\hspace*{2em}\small\textit{Evidence:} ``Neither the HFC nor the HFD evoked a detectable change in plasma TMAO ( \_p\_ > 0.05) in either group.''
\noindent\hspace*{2em}\small\textit{Evidence:} ``There were no significant differences in fasting ( \_p\_ > 0.05) or postprandial TMAO ( \_p\_ > 0.05) between sedentary and endurance‐trained individuals before or following the high‐fat diet, despite the endurance‐trained group consuming more dietary precursors.''
\noindent [79] Constantin-Teodosiu D; Carlin JI; Cederblad G; Harris RC; Hultman E (1991). Acetyl group accumulation and pyruvate dehydrogenase activity in human muscle during incremental exercise.. \textit{Acta Physiol Scand} 143(4):367-72. \href{https://pubmed.ncbi.nlm.nih.gov/1815472/}{PMID:1815472}. \url{https://pubmed.ncbi.nlm.nih.gov/1815472} (accessed 2026-11-15).
\noindent\hspace*{2em}\small\textit{Evidence:} ``The increase of acetyl-CoA was from 12.5 mumol kg-1 dry weight at rest to 27.3 after the highest work load and for acetylcarnitine from 6.0 mmol kg-1 dry weight to 15.2''
\noindent [85] Wall BT; Machin D; Dunlop MV; Stephens FB (2023). Caffeine ingestion stimulates plasma carnitine clearance in humans.. \textit{Physiol Rep} 11(4):e15615. \href{https://doi.org/10.14814/phy2.15615}{doi:10.14814/phy2.15615}. \url{https://pubmed.ncbi.nlm.nih.gov/36806708} (accessed 2026-11-15).
\noindent\hspace*{2em}\small\textit{Evidence:} ``Ingestion of high-dose caffeine stimulates plasma carnitine clearance during hypercarnitinemia, likely via increased Na+ /K+ ATPase activity''
\noindent [86] Longo N; Frigeni M; Pasquali M (2016). Carnitine transport and fatty acid oxidation.. \textit{Biochim Biophys Acta} 1863(10):2422-35. \href{https://doi.org/10.1016/j.bbamcr.2016.01.023}{doi:10.1016/j.bbamcr.2016.01.023}. \url{https://pubmed.ncbi.nlm.nih.gov/26828774} (accessed 2026-11-15).
\noindent\hspace*{2em}\small\textit{Evidence:} ``Defects in the OCTN2 carnitine transporter results in autosomal recessive primary carnitine deficiency characterized by decreased intracellular carnitine accumulation, increased losses of carnitine in the urine, and low serum carnitine levels''
\noindent [88] Seiler SE; Koves TR; Gooding JR; Wong KE; Stevens RD; Ilkayeva OR; Wittmann AH; DeBalsi KL; Davies MN; Lindeboom L; Schrauwen P; Schrauwen-Hinderling VB; Muoio DM (2015). Carnitine Acetyltransferase Mitigates Metabolic Inertia and Muscle Fatigue during Exercise.. \textit{Cell Metab} 22(1):65-76. \href{https://doi.org/10.1016/j.cmet.2015.06.003}{doi:10.1016/j.cmet.2015.06.003}. \url{https://pubmed.ncbi.nlm.nih.gov/26154055} (accessed 2026-11-15).
\noindent\hspace*{2em}\small\textit{Evidence:} ``acetylcarnitine serves as a critical acetyl buffer for working muscles''
\noindent [89] Fujiya M; Inaba Y; Musch MW; Hu S; Kohgo Y; Chang EB (2011). Cytokine regulation of OCTN2 expression and activity in small and large intestine.. \textit{Inflamm Bowel Dis} 17(4):907-16. \href{https://doi.org/10.1002/ibd.21444}{doi:10.1002/ibd.21444}. \url{https://pubmed.ncbi.nlm.nih.gov/20722056} (accessed 2026-11-15).
\noindent\hspace*{2em}\small\textit{Evidence:} ``Treatment with IFN-γ and TNF-α increased intestinal OCTN2 expression, particularly in the colon''
\noindent [91] Evans AM; Fornasini G (2003). Pharmacokinetics of L-carnitine.. \textit{Clin Pharmacokinet} 42(11):941-67. \href{https://pubmed.ncbi.nlm.nih.gov/12908852/}{PMID:12908852}. \url{https://pubmed.ncbi.nlm.nih.gov/12908852} (accessed 2026-11-15).
\noindent\hspace*{2em}\small\textit{Evidence:} ``The threshold concentration for tubular reabsorption (above which the fractional reabsorption begins to decline) is about 40-60 micromol/L, which is similar to the endogenous plasma L-carnitine level''
\noindent\hspace*{2em}\small\textit{Evidence:} ``Under baseline conditions, the renal clearance of L-carnitine (1-3 mL/min) is substantially less than glomerular filtration rate (GFR), indicating extensive (98-99\%) tubular reabsorption.''
\noindent [94] Hatanaka Y; Higuchi T; Akiya Y; Horikami T; Tei R; Furukawa T; Takashima H; Tomita H; Abe M (2019). Prevalence of Carnitine Deficiency and Decreased Carnitine Levels in Patients on Hemodialysis.. \textit{Blood Purif} 47 Suppl 2:38-44. \href{https://doi.org/10.1159/000496720}{doi:10.1159/000496720}. \url{https://pubmed.ncbi.nlm.nih.gov/30943487} (accessed 2026-11-15).
\noindent\hspace*{2em}\small\textit{Evidence:} ``The prevalence of carnitine deficiency and that of carnitine insufficiency was 25.3 and 86.7\%, respectively''
\noindent\hspace*{2em}\small\textit{Evidence:} ``The prevalence of carnitine deficiency and that of carnitine insufficiency was 25.3 and 86.7\%, respectively.''
\noindent [96] Almannai M; Alfadhel M; El-Hattab AW (2019). Carnitine Inborn Errors of Metabolism.. \textit{Molecules} 24(18). \href{https://doi.org/10.3390/molecules24183251}{doi:10.3390/molecules24183251}. \url{https://pubmed.ncbi.nlm.nih.gov/31500110} (accessed 2026-11-15).
\noindent\hspace*{2em}\small\textit{Evidence:} ``In non-vegetarians, most of carnitine sources (\textasciitilde{}75\%) are obtained from diet whereas endogenous synthesis accounts for around 25\%''
\noindent\hspace*{2em}\small\textit{Evidence:} ``One of the major functions of carnitine is shuttling long-chain fatty acids across the mitochondrial membrane from the cytosol into the mitochondrial matrix for β-oxidation. This transport is achieved by mitochondrial carnitine-acylcarnitine cycle, which consists of three enzymes: carnitine palmitoyltransferase I (CPT I), carnitine-acylcarnitine translocase (CACT), and carnitine palmitoyltransferase II (CPT II).''
\noindent\hspace*{2em}\small\textit{Evidence:} ``The biosynthetic pathway involves four enzymes: 6-N-trimethyllysine dioxygenase (TMLD), 3-hydroxy-6-N-trimethyllysine aldolase (HTMLA), 4-N-trimethylaminobutyraldehyde dehydrogenase (TMABADH), and γ-butyrobetaine dioxygenase (BBD). OCTN2 (organic cation/carnitine transporter novel type 2) transports carnitine into the cells''
\noindent [98] Stephens FB; Constantin-Teodosiu D; Laithwaite D; Simpson EJ; Greenhaff PL (2006). Insulin stimulates L-carnitine accumulation in human skeletal muscle.. \textit{FASEB J} 20(2):377-9. \href{https://pubmed.ncbi.nlm.nih.gov/16368715/}{PMID:16368715}. \url{https://pubmed.ncbi.nlm.nih.gov/16368715} (accessed 2026-11-15).
\noindent\hspace*{2em}\small\textit{Evidence:} ``Hypercarnitinemia in the presence of a fasting insulin concentration had no effect on either of these parameters''
\noindent\hspace*{2em}\small\textit{Evidence:} ``The combination of hypercarnitinemia (approximately 500 micromoll(-1)) and hyperinsulinemia increased muscle total carnitine (TC) content from 22.0 +/- 0.9 to 24.7 +/- 1.4 mmol(kg dm)(-1) (P<0.05) and was associated with a 2.3 +/- 0.3-fold increase in carnitine transporter protein (OCTN2) mRNA expression (P<0.05)''
\noindent\hspace*{2em}\small\textit{Evidence:} ``The combination of hypercarnitinemia (approximately 500 micromoll(-1)) and hyperinsulinemia increased muscle total carnitine (TC) content from 22.0 +/- 0.9 to 24.7 +/- 1.4 mmol(kg dm)(-1) (P<0.05)''
\noindent [102] Lu Y; Bai Y; Li W; Zhou Z; Lai H; Hu X; Yang T; Wang C; Chen Y; Gan K; Li K; Ma H; Shen L; He D (2025). Mechanical Stimulation Induces Yap Mediated OCTN2 Transcription to Enhance Carnitine Metabolism in Sarcopenia.. \textit{J Cachexia Sarcopenia Muscle} 16(5):e70052. \href{https://doi.org/10.1002/jcsm.70052}{doi:10.1002/jcsm.70052}. \url{https://pubmed.ncbi.nlm.nih.gov/40958655} (accessed 2026-11-15).
\noindent\hspace*{2em}\small\textit{Evidence:} ``mechanical forces enhanced OCTN2 transcription and regulated carnitine metabolic homeostasis through the Yap/Tead4 pathway''
\noindent [104] Nezu J; Tamai I; Oku A; Ohashi R; Yabuuchi H; Hashimoto N; Nikaido H; Sai Y; Koizumi A; Shoji Y; Takada G; Matsuishi T; Yoshino M; Kato H; Ohura T; Tsujimoto G; Hayakawa J; Shimane M; Tsuji A (1999). Primary systemic carnitine deficiency is caused by mutations in a gene encoding sodium ion-dependent carnitine transporter.. \textit{Nat Genet} 21(1):91-4. \href{https://pubmed.ncbi.nlm.nih.gov/9916797/}{PMID:9916797}. \url{https://pubmed.ncbi.nlm.nih.gov/9916797} (accessed 2026-11-15).
\noindent [107] Tuncel AT; Boy N; Morath MA; Hörster F; Mütze U; Kölker S (2018). Organic acidurias in adults: late complications and management.. \textit{J Inherit Metab Dis} 41(5):765-776. \href{https://doi.org/10.1007/s10545-017-0135-2}{doi:10.1007/s10545-017-0135-2}. \url{https://pubmed.ncbi.nlm.nih.gov/29335813} (accessed 2026-11-15).
\noindent\hspace*{2em}\small\textit{Evidence:} ``Biochemically, OADs are characterized by accumulation of mono-, di- and/or tricarboxylic acids ("organic acids") and corresponding coenzyme A, carnitine and/or glycine esters, some of which are considered toxic at high concentrations.''
\noindent\hspace*{2em}\small\textit{Evidence:} ``Dietary treatment with low protein intake or reduced intake of the precursor amino acid(s), carnitine supplementation, cofactor treatment (in responsive patients) and nonadsorbable antibiotics is commonly used for maintenance treatment.''
\noindent [108] Crefcoeur LL; Visser G; Ferdinandusse S; Wijburg FA; Langeveld M; Sjouke B (2022). Clinical characteristics of primary carnitine deficiency: A structured review using a case-by-case approach.. \textit{J Inherit Metab Dis} 45(3):386-405. \href{https://doi.org/10.1002/jimd.12475}{doi:10.1002/jimd.12475}. \url{https://pubmed.ncbi.nlm.nih.gov/34997761} (accessed 2026-11-15).
\noindent\hspace*{2em}\small\textit{Evidence:} ``In these 621 cases, cardiac symptoms (predominantly cardiomyopathy) were the most prevalent (23.8\%)''
\noindent [109] Frigeni M; Balakrishnan B; Yin X; Calderon FRO; Mao R; Pasquali M; Longo N (2017). Functional and molecular studies in primary carnitine deficiency.. \textit{Hum Mutat} 38(12):1684-1699. \href{https://doi.org/10.1002/humu.23315}{doi:10.1002/humu.23315}. \url{https://pubmed.ncbi.nlm.nih.gov/28841266} (accessed 2026-11-15).
\noindent\hspace*{2em}\small\textit{Evidence:} ``Primary carnitine deficiency is caused by a defect in the OCTN2 carnitine transporter encoded by the SLC22A5 gene. It can cause hypoketotic hypoglycemia or cardiomyopathy in children, and sudden death in children and adults.''
\noindent [110] Novakova K; Török M; Panajatovic M; Bouitbir J; Duong FHT; Handschin C; Krähenbühl S (2022). PGC-1α and MEF2 Regulate the Transcription of the Carnitine Transporter OCTN2 Gene in C2C12 Cells and in Mouse Skeletal Muscle.. \textit{Int J Mol Sci} 23(20). \href{https://doi.org/10.3390/ijms232012304}{doi:10.3390/ijms232012304}. \url{https://pubmed.ncbi.nlm.nih.gov/36293168} (accessed 2026-11-15).
\noindent\hspace*{2em}\small\textit{Evidence:} ``PGC-1α overexpression in C2C12 myoblasts increases OCTN2 mRNA expression''
\noindent [111] Okumura A; Numoto S; Iwayama H; Azuma Y; Kurahashi H (2021). Carnitine supplementation prevents carnitine deficiency caused by pivalate-conjugated antibiotics in patients with epilepsy prescribed valproate.. \textit{Epilepsy Behav} 117:107883. \href{https://doi.org/10.1016/j.yebeh.2021.107883}{doi:10.1016/j.yebeh.2021.107883}. \url{https://pubmed.ncbi.nlm.nih.gov/33714931} (accessed 2026-11-15).
\noindent\hspace*{2em}\small\textit{Evidence:} ``Carnitine deficiency due to PCA use was prevented by carnitine supplementation in patients with epilepsy who were taking VPA''
\noindent [112] Pace S; Longo A; Toon S; Rolan P; Evans AM (2000). Pharmacokinetics of propionyl-L-carnitine in humans: evidence for saturable tubular reabsorption.. \textit{Br J Clin Pharmacol} 50(5):441-8. \href{https://pubmed.ncbi.nlm.nih.gov/11069438/}{PMID:11069438}. \url{https://pubmed.ncbi.nlm.nih.gov/11069438} (accessed 2026-11-15).
\noindent\hspace*{2em}\small\textit{Evidence:} ``At a dose of 1 g, PLC was found to have a mean (+/- s.d.) half-life of 1.09 +/- 0.15 h''
\noindent\hspace*{2em}\small\textit{Evidence:} ``endogenous PLC, LC and ALC underwent extensive renal tubular reabsorption, as indicated by renal excretory clearance to GFR ratios of less than 0.1''
\noindent [113] van Vlies N; Ferdinandusse S; Turkenburg M; Wanders RJ; Vaz FM (2007). PPAR alpha-activation results in enhanced carnitine biosynthesis and OCTN2-mediated hepatic carnitine accumulation.. \textit{Biochim Biophys Acta} 1767(9):1134-42. \href{https://pubmed.ncbi.nlm.nih.gov/17692817/}{PMID:17692817}. \url{https://pubmed.ncbi.nlm.nih.gov/17692817} (accessed 2026-11-15).
\noindent\hspace*{2em}\small\textit{Evidence:} ``both enhancement of carnitine biosynthesis (due to increased gamma-butyrobetaine dioxygenase activity), extra-hepatic gamma-butyrobetaine synthesis and increased hepatic carnitine import (OCTN2 expression) contributes to the increased hepatic carnitine levels after fasting''
\noindent [114] Shannon CE; Nixon AV; Greenhaff PL; Stephens FB (2016). Protein ingestion acutely inhibits insulin-stimulated muscle carnitine uptake in healthy young men.. \textit{Am J Clin Nutr} 103(1):276-82. \href{https://doi.org/10.3945/ajcn.115.119826}{doi:10.3945/ajcn.115.119826}. \url{https://pubmed.ncbi.nlm.nih.gov/26675771} (accessed 2026-11-15).
\noindent\hspace*{2em}\small\textit{Evidence:} ``Protein ingestion acutely inhibits insulin-stimulated muscle carnitine uptake in healthy young men''
\noindent [116] Rebouche CJ (2004). Kinetics, pharmacokinetics, and regulation of L-carnitine and acetyl-L-carnitine metabolism.. \textit{Ann N Y Acad Sci} 1033:30-41. \href{https://pubmed.ncbi.nlm.nih.gov/15591001/}{PMID:15591001}. \url{https://pubmed.ncbi.nlm.nih.gov/15591001} (accessed 2026-11-15).
\noindent\hspace*{2em}\small\textit{Evidence:} ``At normal circulating l-carnitine concentrations, renal l-carnitine reabsorption is highly efficient (90-99\% of filtered load; clearance, 1-3 mL/min), but displays saturation kinetics.''
\noindent\hspace*{2em}\small\textit{Evidence:} ``At normal dietary L-carnitine intake, whole-body turnover time in humans is 38-119 h''
\noindent\hspace*{2em}\small\textit{Evidence:} ``Bioavailability of dietary L-carnitine is 54-87\% and is dependent on the amount of L-carnitine in the meal. Absorption of L-carnitine dietary supplements (0.5-6 g) is primarily passive; bioavailability is 14-18\% of dose.''
\noindent [117] Rebouche CJ; Engel AG (1980). Significance of renal gamma-butyrobetaine hydroxylase for carnitine biosynthesis in man.. \textit{J Biol Chem} 255(18):8700-5. \href{https://pubmed.ncbi.nlm.nih.gov/6773946/}{PMID:6773946}. \url{https://pubmed.ncbi.nlm.nih.gov/6773946} (accessed 2026-11-15).
\noindent\hspace*{2em}\small\textit{Evidence:} ``Significant gamma-butyrobetaine, 2-oxoglutarate dioxygenase (EC 1.14.11.1) activity was found in human kidney but this activity was absent in rat kidney''
\noindent [121] Tamai I; Ohashi R; Nezu J; Yabuuchi H; Oku A; Shimane M; Sai Y; Tsuji A (1998). Molecular and functional identification of sodium ion-dependent, high affinity human carnitine transporter OCTN2.. \textit{J Biol Chem} 273(32):20378-82. \href{https://pubmed.ncbi.nlm.nih.gov/9685390/}{PMID:9685390}. \url{https://pubmed.ncbi.nlm.nih.gov/9685390} (accessed 2026-11-15).
\noindent\hspace*{2em}\small\textit{Evidence:} ``OCTN2 is strongly expressed in kidney, skeletal muscle, heart, and placenta in adult humans''
\noindent\hspace*{2em}\small\textit{Evidence:} ``When OCTN2 was expressed in HEK293 cells, uptake of L-[3H]carnitine was strongly enhanced in a sodium-dependent manner with Km value of 4.34 microM''
\noindent [123] Gziut T; Thanacoody R (2025). L-carnitine for valproic acid-induced toxicity.. \textit{Br J Clin Pharmacol} 91(3):636-647. \href{https://doi.org/10.1111/bcp.16233}{doi:10.1111/bcp.16233}. \url{https://pubmed.ncbi.nlm.nih.gov/39261302} (accessed 2026-11-15).
\noindent\hspace*{2em}\small\textit{Evidence:} ``There is evidence supporting the use of L-carnitine in treating hyperammonaemia and hepatotoxicity following chronic therapeutic use and after acute overdose of valproic acid''
\noindent [126] Stephens FB; Marimuthu K; Cheng Y; Patel N; Constantin D; Simpson EJ; Greenhaff PL (2011). Vegetarians have a reduced skeletal muscle carnitine transport capacity.. \textit{Am J Clin Nutr} 94(3):938-44. \href{https://doi.org/10.3945/ajcn.111.012047}{doi:10.3945/ajcn.111.012047}. \url{https://pubmed.ncbi.nlm.nih.gov/21753065} (accessed 2026-11-15).
\noindent\hspace*{2em}\small\textit{Evidence:} ``Ninety-five percent of the body carnitine pool resides in skeletal muscle where it plays a vital role in fuel metabolism''
\noindent\hspace*{2em}\small\textit{Evidence:} ``[organic cation transporter 2 (OCTN2)] messenger RNA and protein expressions were 16\% (P < 0.01), 58\% (P < 0.01), 17\% (P < 0.05), 33\% (P < 0.05), and 37\% (P = 0.09) lower, respectively, in vegetarian volunteers''
\noindent [127] Linus Pauling Institute (Oregon State University). L-Carnitine — Linus Pauling Institute Micronutrient Information Center. \textit{Oregon State University}. \url{https://lpi.oregonstate.edu/mic/dietary-factors/L-carnitine} (accessed 2026-11-15).
\noindent\hspace*{2em}\small\textit{Evidence:} ``a 70 kg (154 1b) person would synthesize between 11 and 34 mg of carnitine per day''
\noindent\hspace*{2em}\small\textit{Evidence:} ``concentrations of free carnitine around 40 to 50 micromoles (µmol)/L, while plasma concentrations of acetyl-L-carnitine (ALCAR; the most abundant carnitine ester) are around 3 to 6 µmol/L''
\noindent\hspace*{2em}\small\textit{Evidence:} ``This enzyme is, however, highly expressed in human liver, testes, and kidney''
\noindent [128] Kaida Y; Taguchi K; Fukami K (2025). Carnitine Deficiency in Chronic Kidney Disease: Pathophysiology, Clinical Implications, and Therapeutic Perspectives.. \textit{Nutrients} 17(13):. \href{https://doi.org/10.3390/nu17132084}{doi:10.3390/nu17132084}. \url{https://pmc.ncbi.nlm.nih.gov/articles/PMC12251438} (accessed 2026-11-15).
\noindent\hspace*{2em}\small\textit{Evidence:} ``Approximately 70–100 mg of carnitine is shown to be removed during an HD session''
\noindent\hspace*{2em}\small\textit{Evidence:} ``Carnitine is predominantly distributed in skeletal muscle and cardiac muscle, which together account for approximately 98\% of the total body storage. The estimated total carnitine pool is around 100 mmol in a 70 kg adult male''
\noindent\hspace*{2em}\small\textit{Evidence:} ``A large amount of carnitine is filtered through the glomerulus and excreted in the urine, but more than 90\% is reabsorbed in kidney tubular cells.''
\noindent [129] Juraszek B; Nałęcz KA (2019). SLC22A5 (OCTN2) Carnitine Transporter-Indispensable for Cell Metabolism, a Jekyll and Hyde of Human Cancer.. \textit{Molecules} 25(1):. \href{https://doi.org/10.3390/molecules25010014}{doi:10.3390/molecules25010014}. \url{https://pmc.ncbi.nlm.nih.gov/articles/PMC6982704} (accessed 2026-11-15).
\noindent\hspace*{2em}\small\textit{Evidence:} ``co-localization of OCTN2 with PDZK1 (Na+/H+ exchange regulatory cofactor NHE-RF3)''
\noindent [131] van Loon LJ; Greenhaff PL; Constantin-Teodosiu D; Saris WH; Wagenmakers AJ (2001). The effects of increasing exercise intensity on muscle fuel utilisation in humans.. \textit{J Physiol} 536(Pt 1):295-304. \href{https://doi.org/10.1111/j.1469-7793.2001.00295.x}{doi:10.1111/j.1469-7793.2001.00295.x}. \url{https://pubmed.ncbi.nlm.nih.gov/11579177} (accessed 2026-11-15).
\noindent\hspace*{2em}\small\textit{Evidence:} ``We conclude that the most likely mechanism for the reduction in fat oxidation during high-intensity exercise is a downregulation of carnitine palmitoyltransferase I, either by this marked decline in free carnitine availability or by a decrease in intracellular pH.''
\noindent\hspace*{2em}\small\textit{Evidence:} ``This decline involved a decrease in the oxidation rate of both plasma FFA and triacylglycerol fat sources (sum of intramuscular plus lipoprotein-derived triacylglycerol), and was accompanied by increases in muscle pyruvate dehydrogenase complex activation and acetylation of the carnitine pool, resulting in a decline in muscle free carnitine concentration.''
\noindent [135] Wall BT; Stephens FB; Constantin-Teodosiu D; Marimuthu K; Macdonald IA; Greenhaff PL (2011). Chronic oral ingestion of L-carnitine and carbohydrate increases muscle carnitine content and alters muscle fuel metabolism during exercise in humans.. \textit{J Physiol} 589(Pt 4):963-73. \href{https://doi.org/10.1113/jphysiol.2010.201343}{doi:10.1113/jphysiol.2010.201343}. \url{https://pubmed.ncbi.nlm.nih.gov/21224234} (accessed 2026-11-15).
\noindent\hspace*{2em}\small\textit{Evidence:} ``volunteers ingested either 80 g of CHO (Control) or 2 g of l-carnitine-l-tartrate and 80 g of CHO (Carnitine) twice daily for 24 weeks in a randomised, double blind manner.''
\noindent\hspace*{2em}\small\textit{Evidence:} ``Muscle TC increased from basal by 21\% in Carnitine (P < 0.05), and was unchanged in Control.''
\noindent\hspace*{2em}\small\textit{Evidence:} ``At 50\% , the Carnitine group utilised 55\% less muscle glycogen compared to Control (P < 0.05) and 31\% less pyruvate dehydrogenase complex (PDC) activation compared to before supplementation (P < 0.05).''
\noindent [136] Stephens FB; Wall BT; Marimuthu K; Shannon CE; Constantin-Teodosiu D; Macdonald IA; Greenhaff PL (2013). Skeletal muscle carnitine loading increases energy expenditure, modulates fuel metabolism gene networks and prevents body fat accumulation in humans.. \textit{J Physiol} 591(18):4655-66. \href{https://doi.org/10.1113/jphysiol.2013.255364}{doi:10.1113/jphysiol.2013.255364}. \url{https://pubmed.ncbi.nlm.nih.gov/23818692} (accessed 2026-11-15).
\noindent\hspace*{2em}\small\textit{Evidence:} ``whereas muscle total carnitine, long-chain acyl-CoA and whole-body energy expenditure did not change over 12 weeks in Control, they increased in Carnitine by 20\%, 200\% and 6\%, respectively (P<0.05).''
\noindent\hspace*{2em}\small\textit{Evidence:} ``body mass and whole-body fat mass (dual-energy X-ray absorptiometry) increased over 12 weeks in Control by 1.9 and 1.8 kg, respectively (P<0.05), but did not change in Carnitine.''
\noindent [138] Stephens FB; Evans CE; Constantin-Teodosiu D; Greenhaff PL (2007). Carbohydrate ingestion augments L-carnitine retention in humans.. \textit{J Appl Physiol (1985)} 102(3):1065-70. \href{https://doi.org/10.1152/japplphysiol.01011.2006}{doi:10.1152/japplphysiol.01011.2006}. \url{https://pubmed.ncbi.nlm.nih.gov/17138832} (accessed 2026-11-15).
\noindent\hspace*{2em}\small\textit{Evidence:} ``Twenty-four-hour urinary TC excretion in the CHO visit was lower than in the Con visit in study A (155.0 +/- 10.7 vs. 212.1 +/- 17.2 mg; P < 0.05).''
\noindent\hspace*{2em}\small\textit{Evidence:} ``The difference in plasma TC concentration in study A and 24-h urinary TC excretion in both studies suggests that insulin augmented the retention of carnitine in the CHO trials.''
\noindent [139] Stephens FB; Constantin-Teodosiu D; Laithwaite D; Simpson EJ; Greenhaff PL (2007). A threshold exists for the stimulatory effect of insulin on plasma L-carnitine clearance in humans.. \textit{Am J Physiol Endocrinol Metab} 292(2):E637-41. \href{https://doi.org/10.1152/ajpendo.00508.2006}{doi:10.1152/ajpendo.00508.2006}. \url{https://pubmed.ncbi.nlm.nih.gov/17047160} (accessed 2026-11-15).
\noindent\hspace*{2em}\small\textit{Evidence:} ``The findings demonstrate that only high circulating serum insulin concentrations (> or =90 mU/l) are capable of stimulating skeletal muscle carnitine accumulation.''
\noindent [140] Barnett C; Costill DL; Vukovich MD; Cole KJ; Goodpaster BH; Trappe SW; Fink WJ (1994). Effect of L-carnitine supplementation on muscle and blood carnitine content and lactate accumulation during high-intensity sprint cycling.. \textit{Int J Sport Nutr} 4(3):280-8. \href{https://doi.org/10.1123/ijsn.4.3.280}{doi:10.1123/ijsn.4.3.280}. \url{https://pubmed.ncbi.nlm.nih.gov/7987362} (accessed 2026-11-15).
\noindent\hspace*{2em}\small\textit{Evidence:} ``Plasma total acid soluble and free carnitine concentrations were significantly higher (p < .05) at all time points following supplementation. L-carnitine supplementation had no significant effect on muscle carnitine content and thus could not alter lactate accumulation during exercise.''
\noindent\hspace*{2em}\small\textit{Evidence:} ``Following CON II, all subjects began dietary supplementation of L-carnitine for a period of 14 days (4 g/day).''
\noindent [141] Brass EP; Hoppel CL; Hiatt WR (1994). Effect of intravenous L-carnitine on carnitine homeostasis and fuel metabolism during exercise in humans.. \textit{Clin Pharmacol Ther} 55(6):681-92. \href{https://doi.org/10.1038/clpt.1994.85}{doi:10.1038/clpt.1994.85}. \url{https://pubmed.ncbi.nlm.nih.gov/8004884} (accessed 2026-11-15).
\noindent\hspace*{2em}\small\textit{Evidence:} ``Thus the skeletal muscle carnitine pool is segregated from dramatic changes in the plasma carnitine pool, and short-term administration of carnitine has no significant effect on fuel metabolism during exercise in humans.''
\noindent [142] Samborowska E; Olek RA (2023). Twenty-Four Weeks of L-Carnitine Combined with Leucine Supplementation Does Not Increase the Muscle Carnitine Content in Healthy Active Subjects.. \textit{Ann Nutr Metab} 79(2):219-227. \href{https://doi.org/10.1159/000529333}{doi:10.1159/000529333}. \url{https://pubmed.ncbi.nlm.nih.gov/36696886} (accessed 2026-11-15).
\noindent\hspace*{2em}\small\textit{Evidence:} ``Plasma FC and TC content increased in LC + L group after 24 weeks of supplementation (p = 0.003 and 0.010, respectively). However, the skeletal muscle FC and TC contents were not affected by the supplementation protocol.''
\noindent [143] Shannon CE; Ghasemi R; Greenhaff PL; Stephens FB (2018). Increasing skeletal muscle carnitine availability does not alter the adaptations to high-intensity interval training.. \textit{Scand J Med Sci Sports} 28(1):107-115. \href{https://doi.org/10.1111/sms.12885}{doi:10.1111/sms.12885}. \url{https://pubmed.ncbi.nlm.nih.gov/28345160} (accessed 2026-11-15).
\noindent\hspace*{2em}\small\textit{Evidence:} ``Augmenting muscle carnitine during 24 weeks of HIIT did not alter this, nor did it enhance muscle metabolic adaptations or performance gains beyond those with HIIT alone.''
\noindent\hspace*{2em}\small\textit{Evidence:} ``Following HIIT, muscle free carnitine was 30\% greater in CARN vs CON at rest and remained 40\% elevated prior to the start of bout 2 (P<.05).''
\noindent [144] Eryiğit H; Güzel Y; Bulut S; Turnagöl HH; Atakan MM (2026). Acute and Chronic L-Carnitine Supplementation Enhances Endurance Performance Without Consistent Effects on Substrate Oxidation in Active Individuals: A Systematic Review.. \textit{Int J Sport Nutr Exerc Metab} 1-13. \href{https://doi.org/10.1123/ijsnem.2026-0077}{doi:10.1123/ijsnem.2026-0077}. \url{https://pubmed.ncbi.nlm.nih.gov/42887880} (accessed 2026-11-15).
\noindent\hspace*{2em}\small\textit{Evidence:} ``Both acute and chronic L-carnitine supplementation were associated with improvements in endurance performance following acute supplementation and chronic supplementation. By contrast, no significant changes in fat oxidation during exercise were reported after either acute (2-4 g ingested 120-180 min before exercise) or chronic supplementation (2 g/day for 2-4 weeks).''
\noindent\hspace*{2em}\small\textit{Evidence:} ``15 (83\%) studies were rated as having some concerns, three (17\%) were classified as having a high risk of bias, and none were judged to be at low risk of bias.''
\noindent [145] Vecchio M; Chiaramonte R; Testa G; Pavone V (2021). Clinical Effects of L-Carnitine Supplementation on Physical Performance in Healthy Subjects, the Key to Success in Rehabilitation: A Systematic Review and Meta-Analysis from the Rehabilitation Point of View.. \textit{J Funct Morphol Kinesiol} 6. \href{https://doi.org/10.3390/jfmk6040093}{doi:10.3390/jfmk6040093}. \url{https://pubmed.ncbi.nlm.nih.gov/34842765} (accessed 2026-11-15).
\noindent\hspace*{2em}\small\textit{Evidence:} ``The meta-analysis did not show any significant differences in serum lactate values at rest and after exercise in healthy subjects who took L-carnitine supplementation (p > 0.05).''
\noindent [146] Devrim-Lanpir A; Velasco LS; Ramirez Lara FG; Ojeda Sanchez A; Kimble R; Zare R; Gunes FE; Knechtle B; Weiss K; Rosemann T; Heinrich K (2025). Acute L-Carnitine Supplementation Does Not Improve CrossFit(®) Performance: A Randomized, Double-Blind, Placebo-Controlled Crossover Study.. \textit{Nutrients} 17. \href{https://doi.org/10.3390/nu17172784}{doi:10.3390/nu17172784}. \url{https://pubmed.ncbi.nlm.nih.gov/40944177} (accessed 2026-11-15).
\noindent\hspace*{2em}\small\textit{Evidence:} ``The results showed that L-carnitine supplementation did not significantly affect the number of repetitions performed (202.4 ± 69.9 vs. 204.5 ± 78.8, p = 0.810) compared to the placebo.''
\noindent\hspace*{2em}\small\textit{Evidence:} ``However, 10\% of participants reported difficulty sleeping after L-carnitine supplementation.''
\noindent [147] Sawicka AK; Renzi G; Olek RA (2020). The bright and the dark sides of L-carnitine supplementation: a systematic review.. \textit{J Int Soc Sports Nutr} 17(1):49. \href{https://doi.org/10.1186/s12970-020-00377-2}{doi:10.1186/s12970-020-00377-2}. \url{https://pubmed.ncbi.nlm.nih.gov/32958033} (accessed 2026-11-15).
\noindent\hspace*{2em}\small\textit{Evidence:} ``LC supplementation, in combination with carbohydrates (CHO) effectively elevated total carnitine content in skeletal muscle.''
\noindent\hspace*{2em}\small\textit{Evidence:} ``LC supplementation was also noted to induce an increase of fasting plasma trimethylamine-N-oxide (TMAO) levels, which was not associated with modification of determined inflammatory nor oxidative stress markers.''
\noindent [148] Brass EP (2004). Carnitine and sports medicine: use or abuse?. \textit{Ann N Y Acad Sci} 1033():67-78. \href{https://doi.org/10.1196/annals.1320.006}{doi:10.1196/annals.1320.006}. \url{https://pubmed.ncbi.nlm.nih.gov/15591004} (accessed 2026-11-15).
\noindent\hspace*{2em}\small\textit{Evidence:} ``Despite this strong foundation and 20 years of research, no compelling evidence exists that carnitine supplementation can improve physical performance in healthy subjects.''
\noindent\hspace*{2em}\small\textit{Evidence:} ``Muscle carnitine content is not easily increased with carnitine supplementation.''
\noindent [149] Ho JY; Kraemer WJ; Volek JS; Fragala MS; Thomas GA; Dunn-Lewis C; Coday M; Häkkinen K; Maresh CM (2010). l-Carnitine l-tartrate supplementation favorably affects biochemical markers of recovery from physical exertion in middle-aged men and women.. \textit{Metabolism} 59(8):1190-9. \href{https://doi.org/10.1016/j.metabol.2009.11.012}{doi:10.1016/j.metabol.2009.11.012}. \url{https://pubmed.ncbi.nlm.nih.gov/20045157} (accessed 2026-11-15).
\noindent\hspace*{2em}\small\textit{Evidence:} ``Two grams of l-carnitine supplementation had positive effects and significantly (P < or = .05) attenuated biochemical markers of purine metabolism (ie, hypoxanthine, xanthine oxidase), free radical formation (malondialdehyde), muscle tissue disruption (myoglobin, creatine kinase), and muscle soreness after physical exertion. However, markers of physical performance (ie, strength, power, get up and go) were not affected by supplementation.''
\noindent\hspace*{2em}\small\textit{Evidence:} ``Two grams of l-carnitine supplementation had positive effects and significantly (P < or = .05) attenuated biochemical markers of purine metabolism (ie, hypoxanthine, xanthine oxidase), free radical formation (malondialdehyde), muscle tissue disruption (myoglobin, creatine kinase), and muscle soreness after physical exertion.''
\noindent [150] Kraemer WJ; Spiering BA; Volek JS; Ratamess NA; Sharman MJ; Rubin MR; French DN; Silvestre R; Hatfield DL; Van Heest JL; Vingren JL; Judelson DA; Deschenes MR; Maresh CM (2006). Androgenic responses to resistance exercise: effects of feeding and L-carnitine.. \textit{Med Sci Sports Exerc} 38(7):1288-96. \href{https://doi.org/10.1249/01.mss.0000227314.85728.35}{doi:10.1249/01.mss.0000227314.85728.35}. \url{https://pubmed.ncbi.nlm.nih.gov/16826026} (accessed 2026-11-15).
\noindent\hspace*{2em}\small\textit{Evidence:} ``these data demonstrated that: 1) feeding after RE increased AR content, which may result in increased testosterone uptake, and thus enhanced luteinizing hormone secretion via feedback mechanisms; and 2) LCLT supplementation upregulated AR content, which may promote recovery from RE.''
\noindent\hspace*{2em}\small\textit{Evidence:} ``LCLT supplementation upregulated (P<0.05) preexercise AR content compared with PL (12.9+/-5.9 vs 11.2+/-4.0 au, respectively).''
\noindent [151] Spiering BA; Kraemer WJ; Vingren JL; Hatfield DL; Fragala MS; Ho JY; Maresh CM; Anderson JM; Volek JS (2007). Responses of criterion variables to different supplemental doses of L-carnitine L-tartrate.. \textit{J Strength Cond Res} 21(1):259-64. \href{https://doi.org/10.1519/00124278-200702000-00046}{doi:10.1519/00124278-200702000-00046}. \url{https://pubmed.ncbi.nlm.nih.gov/17313301} (accessed 2026-11-15).
\noindent\hspace*{2em}\small\textit{Evidence:} ``The 1- and 2-g doses reduced postexercise serum hypoxanthine, serum xanthine oxidase, serum myoglobin, and perceived muscle soreness.''
\noindent [152] Stefan M; Sharp M; Gheith R; Lowery R; Ottinger C; Wilson J; Durkee S; Bellamine A (2021). L-Carnitine Tartrate Supplementation for 5 Weeks Improves Exercise Recovery in Men and Women: A Randomized, Double-Blind, Placebo-Controlled Trial.. \textit{Nutrients} 13. \href{https://doi.org/10.3390/nu13103432}{doi:10.3390/nu13103432}. \url{https://pubmed.ncbi.nlm.nih.gov/34684429} (accessed 2026-11-15).
\noindent\hspace*{2em}\small\textit{Evidence:} ``Compared to placebo, L-carnitine tartrate supplementation was able to improve perceived recovery and soreness (p = 0.021), and lower serum creatine kinase (p = 0.016).''
\noindent\hspace*{2em}\small\textit{Evidence:} ``Interestingly, serum superoxide dismutase levels increased significantly among those supplementing with L-carnitine tartrate.''
\noindent [153] Caballero-García A; Noriega-González DC; Roche E; Drobnic F; Córdova A (2023). Effects of L-Carnitine Intake on Exercise-Induced Muscle Damage and Oxidative Stress: A Narrative Scoping Review.. \textit{Nutrients} 15. \href{https://doi.org/10.3390/nu15112587}{doi:10.3390/nu15112587}. \url{https://pubmed.ncbi.nlm.nih.gov/37299549} (accessed 2026-11-15).
\noindent\hspace*{2em}\small\textit{Evidence:} ``In view of the studies analyzed and considering the role of L-carnitine in muscle bioenergetics and its antioxidant potential, this supplement could help in postexercise recovery. However, further studies are needed to conclusively clarify the mechanisms underlying these protective effects.''
\noindent [154] Romijn JA; Coyle EF; Sidossis LS; Gastaldelli A; Horowitz JF; Endert E; Wolfe RR (1993). Regulation of endogenous fat and carbohydrate metabolism in relation to exercise intensity and duration.. \textit{Am J Physiol} 265(3 Pt 1):E380-91. \href{https://doi.org/10.1152/ajpendo.1993.265.3.E380}{doi:10.1152/ajpendo.1993.265.3.E380}. \url{https://pubmed.ncbi.nlm.nih.gov/8214047} (accessed 2026-11-15).
\noindent\hspace*{2em}\small\textit{Evidence:} ``In contrast, peripheral lipolysis was stimulated maximally at the lowest exercise intensity, and fatty acid release into plasma decreased with increasing exercise intensity.''
\noindent\hspace*{2em}\small\textit{Evidence:} ``During 2 h of exercise at 65\% VO2max plasma-derived substrate oxidation progressively increased over time, whereas muscle glycogen and triglyceride oxidation decreased.''
\noindent [155] Weiss EP; Frech AM; Perez VR (2023). Low Energy Availability and Increased Risk of Relative Energy Deficiency in Sport (RED-S) During a 3767-km Thru-Hike on the Pacific Crest Trail: A Case Study.. \textit{Wilderness Environ Med} 34(4):536-542. \href{https://doi.org/10.1016/j.wem.2023.06.011}{doi:10.1016/j.wem.2023.06.011}. \url{https://pubmed.ncbi.nlm.nih.gov/37586947} (accessed 2026-11-15).
\noindent\hspace*{2em}\small\textit{Evidence:} ``Exercise energy expenditure above rest (heart rate vs indirect calorimetry regression method) was 2834±1518 kcal/d, total energy expenditure was 5702±1323 kcal/d, and energy intake was 4141 kcal/d. Body mass decreased by 9\%, and fat mass (dual-energy X-ray absorptiometry) decreased by 46\%. Energy availability (energy intake minus exercise energy expenditure) was 19.3 kcal/d/kg fat-free mass, indicating low energy availability (defined as <30 kcal/d/kg).''
\noindent [156] Heinbockel TC; Craighead DH (2021). Case studies in physiology: Impact of a long-distance hike on the Pacific Crest Trail on arterial function and body composition in a highly fit young male.. \textit{Physiol Rep} 9(5):e14767. \href{https://doi.org/10.14814/phy2.14767}{doi:10.14814/phy2.14767}. \url{https://pubmed.ncbi.nlm.nih.gov/33661563} (accessed 2026-11-15).
\noindent\hspace*{2em}\small\textit{Evidence:} ``Individuals who hike the entire length of the trail in one season (4-6 months) perform long daily exercise durations while exposed to extreme environmental temperatures, high altitudes, intense solar radiation, and the consumption of calorie-rich, nutrient-poor diets.''
\noindent [157] Wagner DR; Heath EM; Gustavus M (2026). Case study: Body composition and bone changes following a thru-hike of the Pacific Crest Trail by a young female adult.. \textit{Physiol Rep} 14(15):e71058. \href{https://doi.org/10.14814/phy2.71058}{doi:10.14814/phy2.71058}. \url{https://pubmed.ncbi.nlm.nih.gov/42572415} (accessed 2026-11-15).
\noindent\hspace*{2em}\small\textit{Evidence:} ``Her body fat percentage dropped from 33.2\% to 24.9\%, and there was a 52\% decrease in visceral adipose tissue mass and area.''
\noindent\hspace*{2em}\small\textit{Evidence:} ``A decrease in lumbar BMD despite engaging in load-bearing activity was also observed in two recent case study reports of male thru-hikers, suggesting a maladaptive response that warrants further investigation.''
\noindent [158] Goto H; Nakashima H; Mori K; Tanoue K; Ito S; Kearney BM; Kato A; Nakashima M; Imakiire T; Kumagai H; Kinoshita M; Oshima N (2024). l-Carnitine pretreatment ameliorates heat stress-induced acute kidney injury by restoring mitochondrial function of tubular cells.. \textit{Am J Physiol Renal Physiol} 326(3):F338-F351. \href{https://doi.org/10.1152/ajprenal.00196.2023}{doi:10.1152/ajprenal.00196.2023}. \url{https://pubmed.ncbi.nlm.nih.gov/38095023} (accessed 2026-11-15).
\noindent\hspace*{2em}\small\textit{Evidence:} ``We observed an elevation in serum acylcarnitine levels, indicating carnitine insufficiency in patients with severe heat-related illness and/or AKI.''
\noindent\hspace*{2em}\small\textit{Evidence:} ``l-Carnitine pretreatment ameliorated ATP production in murine tubular cell mitochondria and prevented a change in the kidney macrophage population dynamics observed in AKI: a decrease in tissue-resident macrophages, influx of bone marrow-derived macrophages, and change toward proinflammatory M1 polarization.''
\noindent\hspace*{2em}\small\textit{Evidence:} ``To date, there has been little attention given to the role of carnitine in heat-related illness and AKI.''
\noindent [159] El-Sherbiny HR; El-Shalofy AS; Samir H (2022). Exogenous L-carnitine Administration Ameliorates the Adverse Effects of Heat Stress on Testicular Hemodynamics, Echotexture, and Total Antioxidant Capacity in Rams.. \textit{Front Vet Sci} 9():860771. \href{https://doi.org/10.3389/fvets.2022.860771}{doi:10.3389/fvets.2022.860771}. \url{https://pubmed.ncbi.nlm.nih.gov/35464382} (accessed 2026-11-15).
\noindent\hspace*{2em}\small\textit{Evidence:} ``exogenous LC administration ameliorates testicular hemodynamic disruptions, as measured by spectral Doppler ultrasonography, via augmentation of the rams' total antioxidant capacity under HS conditions.''
\noindent [160] Vukovich MD; Costill DL; Fink WJ (1994). Carnitine supplementation: effect on muscle carnitine and glycogen content during exercise.. \textit{Med Sci Sports Exerc} 26(9):1122-9. \href{https://pubmed.ncbi.nlm.nih.gov/7808246/}{PMID:7808246}. \url{https://pubmed.ncbi.nlm.nih.gov/7808246} (accessed 2026-11-15).
\noindent\hspace*{2em}\small\textit{Evidence:} ``Increased availability of serum carnitine does not result in an increase in muscle carnitine content nor does it alter lipid oxidation.''
\noindent [162] Wächter S; Vogt M; Kreis R; Boesch C; Bigler P; Hoppeler H; Krähenbühl S (2002). Long-term administration of L-carnitine to humans: effect on skeletal muscle carnitine content and physical performance.. \textit{Clin Chim Acta} 318(1-2):51-61. \href{https://doi.org/10.1016/s0009-8981(01)00804-x}{doi:10.1016/s0009-8981(01)00804-x}. \url{https://pubmed.ncbi.nlm.nih.gov/11880112} (accessed 2026-11-15).
\noindent\hspace*{2em}\small\textit{Evidence:} ``The total carnitine content in the skeletal muscle was 4.10 +/- 0.82 micromol/g before, 4.79 +/- 1.19 micromol/g immediately after, and 4.19 +/- 0.61 micromol/g wet weight 2 months after the treatment (no significant difference).''
\noindent\hspace*{2em}\small\textit{Evidence:} ``Beneficial effects of the long-term treatment with L-carnitine on the physical performance of healthy adults cannot be explained by an increase in the carnitine muscle stores.''
\noindent [163] Mansour SG; Verma G; Pata RW; Martin TG; Perazella MA; Parikh CR (2017). Kidney Injury and Repair Biomarkers in Marathon Runners.. \textit{American journal of kidney diseases : the official journal of the National Kidney Foundation} 70(2):252-261. \href{https://doi.org/10.1053/j.ajkd.2017.01.045}{doi:10.1053/j.ajkd.2017.01.045}. \url{https://pubmed.ncbi.nlm.nih.gov/28363731} (accessed 2026-11-15).
\noindent\hspace*{2em}\small\textit{Evidence:} ``82\% of runners developed an increase in creatinine level equivalent to AKIN-defined AKI stages 1 and 2. 73\% had microscopy diagnoses of tubular injury.''
\noindent\hspace*{2em}\small\textit{Evidence:} ``An increase in injury and repair biomarker levels suggests structural damage to renal tubules occurring after marathon.''
\noindent\hspace*{2em}\small\textit{Evidence:} ``73\% had microscopy diagnoses of tubular injury.''
\noindent [164] Hoffman MD; Weiss RH (2016). Does Acute Kidney Injury From an Ultramarathon Increase the Risk for Greater Subsequent Injury?. \textit{Clinical journal of sport medicine : official journal of the Canadian Academy of Sport Medicine} 26(5):417-22. \href{https://doi.org/10.1097/JSM.0000000000000277}{doi:10.1097/JSM.0000000000000277}. \url{https://pubmed.ncbi.nlm.nih.gov/26657822} (accessed 2026-11-15).
\noindent\hspace*{2em}\small\textit{Evidence:} ``Among 627 finishes in which serum creatinine values were known, 36.2\% met "risk" or "injury" criterion with this group characterized by having faster finish times, greater body weight loss during the race, and higher postrace serum creatine kinase and urea nitrogen concentrations when compared with those not meeting the criteria.''
\noindent\hspace*{2em}\small\textit{Evidence:} ``Mild AKI is common in 161-km ultramarathons, but there was no evidence that previous AKI caused greater renal dysfunction from a subsequent exercise stimulus of similar magnitude.''
\noindent\hspace*{2em}\small\textit{Evidence:} ``For most (56.2\%) of the 16 runners meeting the criteria at the first race, the subsequent race caused less increase in serum creatinine concentration and decrement in estimated glomerular filtration rate than the first race.''
\noindent [165] Poussel M; Touzé C; Allado E; Frimat L; Hily O; Thilly N; Rousseau H; Vauthier JC; Chenuel B (2019). Ultramarathon and Renal Function: Does Exercise-Induced Acute Kidney Injury Really Exist in Common Conditions?. \textit{Frontiers in sports and active living} 1:71. \href{https://doi.org/10.3389/fspor.2019.00071}{doi:10.3389/fspor.2019.00071}. \url{https://pubmed.ncbi.nlm.nih.gov/33344994} (accessed 2026-11-15).
\noindent\hspace*{2em}\small\textit{Evidence:} ``Depending the methodology used to calculate GFR, the prevalence of AKI was observed from 0 to 12.5\%.''
\noindent\hspace*{2em}\small\textit{Evidence:} ``Our study showed a very low prevalence of AKI and no evidence that ultra-endurance running can cause important kidney damage in properly hydrated subjects with no use of NSAIDs.''
\noindent\hspace*{2em}\small\textit{Evidence:} ``Whether the increase in urinary biomarkers of kidney damage following the race reflects structural kidney injury or a simple metabolic adaptation to strenuous exercise needs to be clarified.''
\noindent [166] Leucuța DC; Trif LI; Almășan O; Popa ȘL; Ismaiel A (2025). Acute Kidney Injury Biomarkers in Marathon Runners: Systematic Review and Meta-Analysis.. \textit{Medicina (Kaunas, Lithuania)} 61(10):. \href{https://doi.org/10.3390/medicina61101775}{doi:10.3390/medicina61101775}. \url{https://pubmed.ncbi.nlm.nih.gov/41155766} (accessed 2026-11-15).
\noindent\hspace*{2em}\small\textit{Evidence:} ``Marathon running causes significant increases in kidney injury biomarkers, with different patterns of evolution.''
\noindent\hspace*{2em}\small\textit{Evidence:} ``The largest increases were observed in Tissue Inhibitor of Metalloproteinases-2 Insulin-like Growth Factor Binding Protein-7 (TIMP-2IGFBP), copeptin, urinary Liver-type Fatty Acid Binding Protein (L-FABP), urinary Monocyte Chemoattractant Protein-1 (MCP-1), IGFBP-7, urinary Chitinase 3-like Prote''
\noindent [167] Lipman GS; Krabak BJ; Rundell SD; Shea KM; Badowski N; Little C (2016). Incidence and Prevalence of Acute Kidney Injury During Multistage Ultramarathons.. \textit{Clinical journal of sport medicine : official journal of the Canadian Academy of Sport Medicine} 26(4):314-9. \href{https://doi.org/10.1097/JSM.0000000000000253}{doi:10.1097/JSM.0000000000000253}. \url{https://pubmed.ncbi.nlm.nih.gov/26513390} (accessed 2026-11-15).
\noindent\hspace*{2em}\small\textit{Evidence:} ``Stage 1 had 56 (43.8\%) with risk of AKI and 24 (18.8\%) with injury; in stage 3, 61 (47.7\%) were at risk, 41 (32\%) had injury; in stage 5, 62 (48.4\%) runners were at risk and 36 (28.1\%) had injury.''
\noindent\hspace*{2em}\small\textit{Evidence:} ``Cumulative incidence of AKI was 41.4\%.''
\noindent\hspace*{2em}\small\textit{Evidence:} ``Female sex, lower pack weight, and greater weight loss were associated with renal impairment.''
\noindent [168] Lipman GS; Krabak BJ; Waite BL; Logan SB; Menon A; Chan GK (2014). A prospective cohort study of acute kidney injury in multi-stage ultramarathon runners: the Biochemistry in Endurance Runner Study (BIERS).. \textit{Research in sports medicine (Print)} 22(2):185-92. \href{https://doi.org/10.1080/15438627.2014.881824}{doi:10.1080/15438627.2014.881824}. \url{https://pubmed.ncbi.nlm.nih.gov/24650338} (accessed 2026-11-15).
\noindent\hspace*{2em}\small\textit{Evidence:} ``no apparent cumulative effect was observed, with resolution of renal function to near baseline levels between stages.''
\noindent\hspace*{2em}\small\textit{Evidence:} ``the majority of participants (55-80\%) incurring AKI.''
\noindent [169] Chapman CL; Johnson BD; Parker MD; Hostler D; Pryor RR; Schlader Z (2021). Kidney physiology and pathophysiology during heat stress and the modification by exercise, dehydration, heat acclimation and aging.. \textit{Temperature (Austin, Tex.)} 8(2):108-159. \href{https://doi.org/10.1080/23328940.2020.1826841}{doi:10.1080/23328940.2020.1826841}. \url{https://pubmed.ncbi.nlm.nih.gov/33997113} (accessed 2026-11-15).
\noindent\hspace*{2em}\small\textit{Evidence:} ``Recent evidence suggests the kidneys are at increased risk of pathological events during heat stress, namely acute kidney injury (AKI), and that this risk is compounded by dehydration and exercise.''
\noindent\hspace*{2em}\small\textit{Evidence:} ``This heat stress related AKI is believed to contribute to the epidemic of chronic kidney disease (CKD) occurring in occupational settings.''
\noindent [170] García-Trabanino R; Jarquín E; Wesseling C; Johnson RJ; González-Quiroz M; Weiss I; Glaser J; José Vindell J; Stockfelt L; Roncal C; Harra T; Barregard L (2015). Heat stress, dehydration, and kidney function in sugarcane cutters in El Salvador--A cross-shift study of workers at risk of Mesoamerican nephropathy.. \textit{Environmental research} 142:746-55. \href{https://doi.org/10.1016/j.envres.2015.07.007}{doi:10.1016/j.envres.2015.07.007}. \url{https://pubmed.ncbi.nlm.nih.gov/26209462} (accessed 2026-11-15).
\noindent\hspace*{2em}\small\textit{Evidence:} ``Serum creatinine, uric acid and urea nitrogen increased, while chloride and potassium decreased.''
\noindent\hspace*{2em}\small\textit{Evidence:} ``Mean workday temperature was 34-36 °C before noon, and 39-42 °C at noon.''
\noindent\hspace*{2em}\small\textit{Evidence:} ``The high prevalence of reduced eGFR, and the cross-shift changes are consistent with recurrent dehydration from strenuous work in a hot and humid environment as an important causal factor.''
\noindent [171] Roncal-Jimenez C; García-Trabanino R; Barregard L; Lanaspa MA; Wesseling C; Harra T; Aragón A; Grases F; Jarquin ER; González MA; Weiss I; Glaser J; Sánchez-Lozada LG; Johnson RJ (2016). Heat Stress Nephropathy From Exercise-Induced Uric Acid Crystalluria: A Perspective on Mesoamerican Nephropathy.. \textit{American journal of kidney diseases : the official journal of the National Kidney Foundation} 67(1):20-30. \href{https://doi.org/10.1053/j.ajkd.2015.08.021}{doi:10.1053/j.ajkd.2015.08.021}. \url{https://pubmed.ncbi.nlm.nih.gov/26455995} (accessed 2026-11-15).
\noindent\hspace*{2em}\small\textit{Evidence:} ``Increased hydration with water and salt, urinary alkalinization, reduction in sugary beverage intake, and inhibitors of uric acid synthesis should be tested for disease prevention.''
\noindent\hspace*{2em}\small\textit{Evidence:} ``Individuals at risk for developing the disease are primarily male workers exposed to heat stress and physical exertion that predisposes to recurrent water and volume depletion, often accompanied by urinary concentration and acidification.''
\noindent\hspace*{2em}\small\textit{Evidence:} ``Uric acid is generated during heat stress, in part consequent to nucleotide release from muscles.''
\noindent [172] Amorim F; Schlader Z (2025). The kidney under heat stress: a vulnerable state.. \textit{Current opinion in nephrology and hypertension} 34(2):170-176. \href{https://doi.org/10.1097/MNH.0000000000001050}{doi:10.1097/MNH.0000000000001050}. \url{https://pubmed.ncbi.nlm.nih.gov/39688252} (accessed 2026-11-15).
\noindent\hspace*{2em}\small\textit{Evidence:} ``As body temperature increases, blood flow is directed toward the skin to aid in cooling, diverting it away from internal organs like the kidneys to support blood pressure regulation.''
\noindent\hspace*{2em}\small\textit{Evidence:} ``Collectively, this can create a localized supply-demand mismatch, resulting in tissue hypoxia that can damage kidney tissues.''
\noindent [173] DeGroot DW; Goodwin KC; Giersch GEW; Charkoudian N (2026). Acute kidney injury and liver damage are the most common forms of end organ damage in exertional heat illness in otherwise healthy young adults.. \textit{Temperature (Austin, Tex.)} 13(3):264-274. \href{https://doi.org/10.1080/23328940.2026.2680846}{doi:10.1080/23328940.2026.2680846}. \url{https://pubmed.ncbi.nlm.nih.gov/42694931} (accessed 2026-11-15).
\noindent\hspace*{2em}\small\textit{Evidence:} ``Acute kidney injury was the most common form of end organ damage observed (73-86\%) followed closely by liver damage (59-73\%).''
\noindent\hspace*{2em}\small\textit{Evidence:} ``EHS cases had higher peak damage levels for Cr, CK, AST, and ALT compared to either HI or HE.''
\noindent [174] Melin B; Jimenez C; Savourey G; Bittel J; Cottet-Emard JM; Pequignot JM; Allevard AM; Gharib C (1997). Effects of hydration state on hormonal and renal responses during moderate exercise in the heat.. \textit{European journal of applied physiology and occupational physiology} 76(4):320-7. \href{https://doi.org/10.1007/s004210050255}{doi:10.1007/s004210050255}. \url{https://pubmed.ncbi.nlm.nih.gov/9349646} (accessed 2026-11-15).
\noindent\hspace*{2em}\small\textit{Evidence:} ``During exercise in the Dh compared to C state, plasma renin, aldosterone, arginine vasopressin (AVP), noradrenaline and adrenaline concentrations were increased (P < 0.05).''
\noindent\hspace*{2em}\small\textit{Evidence:} ``A reduction in creatinine clearance and urine flow was also observed (P < 0.05) together with a decrease in''
\noindent\hspace*{2em}\small\textit{Evidence:} ``Rehydration, by reducing this activation, could partially restore the renal concentrating ability despite the lowered plasma AVP.''
\noindent [175] Becerra LA; Mansour SG (2026). Exercise and Kidney Health: Core Curriculum 2026.. \textit{American journal of kidney diseases : the official journal of the National Kidney Foundation} 87(2):246-259. \href{https://doi.org/10.1053/j.ajkd.2025.09.014}{doi:10.1053/j.ajkd.2025.09.014}. \url{https://pubmed.ncbi.nlm.nih.gov/41389078} (accessed 2026-11-15).
\noindent\hspace*{2em}\small\textit{Evidence:} ``Prevention strategies include individualized hydration plans, electrolyte replacement, and avoidance of nonsteroidal anti-inflammatory drugs.''
\noindent\hspace*{2em}\small\textit{Evidence:} ``Adaptations in kidney physiology during exercise include reduced plasma flow, altered glomerular filtration, and hormone-mediated fluid retention.''
\noindent\hspace*{2em}\small\textit{Evidence:} ``These changes are protective but may become maladaptive with dehydration, heat stress, or excessive fluid intake.''
\noindent [176] Rojas-Valverde D; Sánchez-Ureña B; Crowe J; Timón R; Olcina GJ (2021). Exertional rhabdomyolysis and acute kidney injury in endurance sports: A systematic review.. \textit{European journal of sport science} 21(2):261-274. \href{https://doi.org/10.1080/17461391.2020.1746837}{doi:10.1080/17461391.2020.1746837}. \url{https://pubmed.ncbi.nlm.nih.gov/32202487} (accessed 2026-11-15).
\noindent\hspace*{2em}\small\textit{Evidence:} ``345 (43.5\%) individuals were diagnosed with ER (creatinine kinase > 5000 UI/L) and 130 (16.39\%) with ER + AKI (creatinine ≥ 1.88 mg/dL).''
\noindent\hspace*{2em}\small\textit{Evidence:} ``Out of the total cases of ER + AKI, 96.92\% were in ultra-endurance runners.''
\noindent\hspace*{2em}\small\textit{Evidence:} ``Increased levels of muscle and kidney injury immediately after endurance events were reported, but after 5.86 days these levels usually returned to baseline.''
\noindent [177] Sabouri AH; Yurgionas B; Khorasani S; Durant EJ; Kafaie J; Hung YY; Klingman JG; Nemazie S (2024). Acute Kidney Injury in Hospitalized Patients With Exertional Rhabdomyolysis.. \textit{JAMA network open} 7(8):e2427464. \href{https://doi.org/10.1001/jamanetworkopen.2024.27464}{doi:10.1001/jamanetworkopen.2024.27464}. \url{https://pubmed.ncbi.nlm.nih.gov/39136944} (accessed 2026-11-15).
\noindent\hspace*{2em}\small\textit{Evidence:} ``the risk of AKI was significantly higher in patients with ERM who used nonsteroidal anti-inflammatory drugs (NSAIDs) before admission (11 of 17 with AKI [64.7\%] vs 40 of 183 without AKI [21.9\%], P < .001) or experienced dehydration (9 of 183 without AKI [52.9\%] vs 9 of 17 with AKI [4.9\%], P < .001).''
\noindent\hspace*{2em}\small\textit{Evidence:} ``Seventeen patients (8.5\%) developed AKI, none developed CKD, 1 (0.5\%) developed compartment syndrome, and there were no fatalities.''
\noindent\hspace*{2em}\small\textit{Evidence:} ``There was no association between serum CK levels and the risk of AKI.''
\noindent [179] Myers TM; Bennett BL; Myers W (2024). Exercise-Associated Hyponatremic Encephalopathy, Rhabdomyolysis, and Acute Kidney Injury While Hiking in the Grand Canyon.. \textit{Wilderness \& environmental medicine} 35(4):484-489. \href{https://doi.org/10.1177/10806032241286487}{doi:10.1177/10806032241286487}. \url{https://pubmed.ncbi.nlm.nih.gov/39444363} (accessed 2026-11-15).
\noindent\hspace*{2em}\small\textit{Evidence:} ``We present an unusual case of concomitant exercise-associated hyponatremic encephalopathy (EAHE), exertional rhabdomyolysis (ER), and acute kidney injury (AKI) in a Grand Canyon hiker.''
\noindent\hspace*{2em}\small\textit{Evidence:} ``consumed an excessive amount of water during his descent into the Canyon during hot weather.''
\noindent\hspace*{2em}\small\textit{Evidence:} ``This case demonstrates the seriousness of overconsumption of water, as well as the potential complications of muscle damage when a deconditioned person does prolonged, strenuous exercise.''
\noindent [180] Maeo S; Yamamoto M; Kanehisa H; Nosaka K (2017). Prevention of downhill walking-induced muscle damage by non-damaging downhill walking.. \textit{PloS one} 12(3):e0173909. \href{https://doi.org/10.1371/journal.pone.0173909}{doi:10.1371/journal.pone.0173909}. \url{https://pubmed.ncbi.nlm.nih.gov/28288187} (accessed 2026-11-15).
\noindent\hspace*{2em}\small\textit{Evidence:} ``These changes were significantly (P<0.05) attenuated (47-64\%) for the PRE-DW group (-9.9 ± 9.6\%, 339.3 ± 148.4 U/L, 27.8 ± 16.8 mm).''
\noindent\hspace*{2em}\small\textit{Evidence:} ``Prolonged DW induces damage to leg muscles, reducing force generating ability and muscle coordination.''
\noindent\hspace*{2em}\small\textit{Evidence:} ``40-min DW (-28\%, 5 km/h, 10\% load).''
\noindent [181] Lipman GS; Shea K; Christensen M; Phillips C; Burns P; Higbee R; Koskenoja V; Eifling K; Krabak BJ (2017). Ibuprofen versus placebo effect on acute kidney injury in ultramarathons: a randomised controlled trial.. \textit{Emergency medicine journal : EMJ} 34(10):637-642. \href{https://doi.org/10.1136/emermed-2016-206353}{doi:10.1136/emermed-2016-206353}. \url{https://pubmed.ncbi.nlm.nih.gov/28679502} (accessed 2026-11-15).
\noindent\hspace*{2em}\small\textit{Evidence:} ``Intent-to-treat analysis found 22 (52\%) ibuprofen versus 16 (34\%) placebo users developed AKI (18\% difference, 95\% CI -4\% to 41\%; OR 2.1, 95\% CI 0.9 to 5.1) with a number needed to harm of 5.5.''
\noindent\hspace*{2em}\small\textit{Evidence:} ``up to 75\% of ultramarathon runners ingest these during competition''
\noindent\hspace*{2em}\small\textit{Evidence:} ``The overall incidence of AKI was 44\%.''
\noindent [182] Cole PL; Bontemps B; Pugh J; Owens DJ (2026). The prevalence and perceptions of non-steroidal anti-inflammatory drug use among ultra-endurance runners.. \textit{Research in sports medicine (Print)} 34(3):257-275. \href{https://doi.org/10.1080/15438627.2025.2561040}{doi:10.1080/15438627.2025.2561040}. \url{https://pubmed.ncbi.nlm.nih.gov/40947642} (accessed 2026-11-15).
\noindent\hspace*{2em}\small\textit{Evidence:} ``Over half (53\%) reported NSAID use, with ibuprofen most common.''
\noindent\hspace*{2em}\small\textit{Evidence:} ``This mixed-methods, cross-sectional survey included 167 ultra-endurance runners across five self-reported performance tiers.''
\noindent\hspace*{2em}\small\textit{Evidence:} ``Despite 77\% being aware of potential harms, most often renal and gastrointestinal (GI), usage remained high due to perceived benefits.''
\noindent [183] McDermott BP; Smith CR; Butts CL; Caldwell AR; Lee EC; Vingren JL; Munoz CX; Kunces LJ; Williamson K; Ganio MS; Armstrong LE (2018). Renal stress and kidney injury biomarkers in response to endurance cycling in the heat with and without ibuprofen.. \textit{Journal of science and medicine in sport} 21(12):1180-1184. \href{https://doi.org/10.1016/j.jsams.2018.05.003}{doi:10.1016/j.jsams.2018.05.003}. \url{https://pubmed.ncbi.nlm.nih.gov/29784554} (accessed 2026-11-15).
\noindent\hspace*{2em}\small\textit{Evidence:} ``Changes in renal biomarkers suggest mild acute kidney injury and reduced kidney function during a single bout of endurance cycling in the heat, without influence from moderate ibuprofen ingestion.''
\noindent\hspace*{2em}\small\textit{Evidence:} ``Serum creatinine significantly increased from pre- (0.52±0.14mg/dL) to post-event (0.88±0.21mg/dL; p<0.001).''
\noindent [184] Specht JW; Bailly AR; Garcia S; Klepacz S; Andrade DE Oliveira S; Lucero D; McKenna ZJ; Schlader ZJ; Amorim FT (2025). Effect of Ibuprofen on Markers of Acute Kidney Injury, Intestinal Injury, and Endotoxemia after Running in the Heat.. \textit{Medicine and science in sports and exercise} 57(6):1092-1102. \href{https://doi.org/10.1249/MSS.0000000000003659}{doi:10.1249/MSS.0000000000003659}. \url{https://pubmed.ncbi.nlm.nih.gov/39876077} (accessed 2026-11-15).
\noindent\hspace*{2em}\small\textit{Evidence:} ``These findings indicate that ibuprofen ingestion does not worsen intestinal or renal injury experienced during 1 h of exercise in the heat, but increases pro-inflammatory IL-8.''
\noindent [187] Crouse BJ; Josephs D (1993). Health Care Needs of Appalachian Trail Hikers. \textit{The Journal of Family Practice} 36(5):521-525. \href{https://pubmed.ncbi.nlm.nih.gov/8482936/}{PMID:8482936}. \url{https://cdn.mdedge.com/files/s3fs-public/jfp-archived-issues/1993-volume_36-37/JFP_1993-05_v36_i5_health-care-needs-of-appalachian-trail-h.pdf} (accessed 2026-11-15).
\noindent\hspace*{2em}\small\textit{Evidence:} ``Analgesics and nonsteroidal anti-inflammatory agents were frequently carried and used.''
\noindent\hspace*{2em}\small\textit{Evidence:} ``The majority of the analgesic medications were nonsteroidal, anti-inflammatory agents.''
\noindent\hspace*{2em}\small\textit{Evidence:} ``Injuries and illnesses were reported by 82\% of the respondents, incurring an average loss of 4.7 days of hiking.''
\noindent [188] Spano SJ; Hile AG; Jain R; Stalcup PR (2018). The Epidemiology and Medical Morbidity of Long-Distance Backpackers on the John Muir Trail in the Sierra Nevada.. \textit{Wilderness \& environmental medicine} 29(2):203-210. \href{https://doi.org/10.1016/j.wem.2018.02.006}{doi:10.1016/j.wem.2018.02.006}. \url{https://pubmed.ncbi.nlm.nih.gov/29804621} (accessed 2026-11-15).
\noindent\hspace*{2em}\small\textit{Evidence:} ``Backpackers lost 3.5±2.6 (+3.6 to -18.2) kg (7.7±5.8 [+8 to -40] lbs).''
\noindent\hspace*{2em}\small\textit{Evidence:} ``Over half (57\%) of respondents reported illness or injury, with blisters (57\%), sleep problems (57\%), and pack strap pain (46\%) most prevalent.''
\noindent [189] Almond CS; Shin AY; Fortescue EB; Mannix RC; Wypij D; Binstadt BA; Duncan CN; Olson DP; Salerno AE; Newburger JW; Greenes DS (2005). Hyponatremia among runners in the Boston Marathon.. \textit{The New England journal of medicine} 352(15):1550-6. \href{https://doi.org/10.1056/NEJMoa043901}{doi:10.1056/NEJMoa043901}. \url{https://pubmed.ncbi.nlm.nih.gov/15829535} (accessed 2026-11-15).
\noindent\hspace*{2em}\small\textit{Evidence:} ``On multivariate analysis, hyponatremia was associated with weight gain (odds ratio, 4.2; 95 percent confidence interval, 2.2 to 8.2), a racing time of >4:00 hours (odds ratio for the comparison with a time of <3:30 hours, 7.4; 95 percent confidence interval, 2.9 to 23.1), and body-mass-index extremes.''
\noindent\hspace*{2em}\small\textit{Evidence:} ``Thirteen percent had hyponatremia (a serum sodium concentration of 135 mmol per liter or less); 0.6 percent had critical hyponatremia (120 mmol per liter or less).''
\noindent [190] Bennett BL; Hew-Butler T; Rosner MH; Myers T; Lipman GS (2020). Wilderness Medical Society Clinical Practice Guidelines for the Management of Exercise-Associated Hyponatremia: 2019 Update.. \textit{Wilderness \& environmental medicine} 31(1):50-62. \href{https://doi.org/10.1016/j.wem.2019.11.003}{doi:10.1016/j.wem.2019.11.003}. \url{https://pubmed.ncbi.nlm.nih.gov/32044213} (accessed 2026-11-15).
\noindent\hspace*{2em}\small\textit{Evidence:} ``Exercise-associated hyponatremia (EAH) is defined by a serum or plasma sodium concentration below the normal reference range of 135 mmol·L-1 that occurs during or up to 24 h after prolonged physical activity.''
\noindent\hspace*{2em}\small\textit{Evidence:} ``EAH most commonly has overlapping signs and symptoms with heat exhaustion and exertional heat stroke.''
\noindent [192] Noe RS; Choudhary E; Cheng-Dobson J; Wolkin AF; Newman SB (2013). Exertional heat-related illnesses at the Grand Canyon National Park, 2004-2009.. \textit{Wilderness \& environmental medicine} 24(4):422-8. \href{https://doi.org/10.1016/j.wem.2013.06.008}{doi:10.1016/j.wem.2013.06.008}. \url{https://pubmed.ncbi.nlm.nih.gov/24119571} (accessed 2026-11-15).
\noindent\hspace*{2em}\small\textit{Evidence:} ``Clinical HRI assessments included dehydration (25\%), heat exhaustion (23\%), and suspected hyponatremia (19\%).''
\noindent\hspace*{2em}\small\textit{Evidence:} ``Grand Canyon EMS responded to 474 nonfatal and 6 fatal HRI cases, with the majority''
\noindent\hspace*{2em}\small\textit{Evidence:} ``Almost all (90\%) were hiking; 40\% required helicopter evacuation.''
\noindent [194] Connor KR; Brazelton SC; Greenfield AM; Charkoudian N; Chapman CL; Giersch GEW (2025). Heat acclimation alters the relation between hyperthermia and biomarkers of kidney function during exercise heat stress.. \textit{American journal of physiology. Regulatory, integrative and comparative physiology} 329(6):R995-R1001. \href{https://doi.org/10.1152/ajpregu.00215.2025}{doi:10.1152/ajpregu.00215.2025}. \url{https://pubmed.ncbi.nlm.nih.gov/41247794} (accessed 2026-11-15).
\noindent\hspace*{2em}\small\textit{Evidence:} ``Heat acclimation did not alter the increase in ΔCr and ΔNGAL to exercise heat stress (P ≥ 0.325).''
\noindent\hspace*{2em}\small\textit{Evidence:} ``the kidneys may exhibit increased sensitivity to hyperthermia in heat-acclimated individuals.NEW \& NOTEWORTHY He''
\noindent [195] Omassoli J; Hill NE; Woods DR; Delves SK; Fallowfield JL; Brett SJ; Wilson D; Corbett RW; Allsopp AJ; Stacey MJ (2019). Variation in renal responses to exercise in the heat with progressive acclimatisation.. \textit{Journal of science and medicine in sport} 22(9):1004-1009. \href{https://doi.org/10.1016/j.jsams.2019.04.010}{doi:10.1016/j.jsams.2019.04.010}. \url{https://pubmed.ncbi.nlm.nih.gov/31085130} (accessed 2026-11-15).
\noindent\hspace*{2em}\small\textit{Evidence:} ``Heat acclimatization is associated with reduced markers of renal stress and AKI incidence, perhaps due to improved regional perfusion.''
\noindent\hspace*{2em}\small\textit{Evidence:} ``Acute Kidney Injury was common (26/100 exposures) and occurred most frequently in the unacclimatised state.''
\noindent\hspace*{2em}\small\textit{Evidence:} ``Acclimatisation and monitoring Ratings of Perceived Exertion are practical, non-invasive measures that could help to reduce renal injury from exercise in the heat.''
\noindent [196] Nishioka N; Luo Y; Taniguchi T; Ohnishi T; Kimachi M; Ng RC; Watanabe N (2022). Carnitine supplements for people with chronic kidney disease requiring dialysis.. \textit{Cochrane Database Syst Rev} 12(12):CD013601. \href{https://doi.org/10.1002/14651858.CD013601.pub2}{doi:10.1002/14651858.CD013601.pub2}. \url{https://pubmed.ncbi.nlm.nih.gov/36472884} (accessed 2026-11-15).
\noindent\hspace*{2em}\small\textit{Evidence:} ``The available evidence does not currently support the use of carnitine supplementation in the treatment of dialysis-related carnitine deficiency.''
\noindent\hspace*{2em}\small\textit{Evidence:} ``Although carnitine supplementation may slightly improve anaemia-related markers, carnitine supplementation makes little or no difference to adverse events.''
\noindent [197] Chen Y; Abbate M; Tang L; Cai G; Gong Z; Wei R; Zhou J; Chen X (2014). L-Carnitine supplementation for adults with end-stage kidney disease requiring maintenance hemodialysis: a systematic review and meta-analysis.. \textit{Am J Clin Nutr} 99(2):408-22. \href{https://doi.org/10.3945/ajcn.113.062802}{doi:10.3945/ajcn.113.062802}. \url{https://pubmed.ncbi.nlm.nih.gov/24368434} (accessed 2026-11-15).
\noindent\hspace*{2em}\small\textit{Evidence:} ``This meta-analysis failed to confirm the previous findings regarding the effects of l-carnitine on hemoglobin and the erythropoietin dose but showed that l-carnitine significantly decreased serum LDL and CRP.''
\noindent [198] Higuchi T; Abe M; Yamazaki T; Okawa E; Ando H; Hotta S; Oikawa O; Kikuchi F; Okada K; Soma M (2016). Levocarnitine Improves Cardiac Function in Hemodialysis Patients With Left Ventricular Hypertrophy: A Randomized Controlled Trial.. \textit{Am J Kidney Dis} 67(2):260-70. \href{https://doi.org/10.1053/j.ajkd.2015.09.010}{doi:10.1053/j.ajkd.2015.09.010}. \url{https://pubmed.ncbi.nlm.nih.gov/26508680} (accessed 2026-11-15).
\noindent\hspace*{2em}\small\textit{Evidence:} ``Levocarnitine therapy is useful for hemodialysis patients with carnitine deficiency; these patients may benefit from such therapy, with amelioration of cardiac function and reduction of left ventricular mass index.''
\noindent [199] Reuter SE; Steiber AL; Faull RJ; Evans AM (2026). Effectiveness of L-Carnitine for the Treatment of Erythropoietin-Resistant Renal Anemia: A Randomized, Double-Blind, Placebo-Controlled Pilot Trial.. \textit{J Ren Nutr} 36(3):491-501. \href{https://doi.org/10.1053/j.jrn.2026.01.008}{doi:10.1053/j.jrn.2026.01.008}. \url{https://pubmed.ncbi.nlm.nih.gov/41644022} (accessed 2026-11-15).
\noindent\hspace*{2em}\small\textit{Evidence:} ``L-carnitine supplementation was associated with a 25\% reduction in ERI, significantly greater than that seen in placebo-treated patients (3\% reduction).''
\noindent [200] Evans A (2003). Dialysis-related carnitine disorder and levocarnitine pharmacology.. \textit{Am J Kidney Dis} 41(4 Suppl 4):S13-26. \href{https://doi.org/10.1016/s0272-6386(03)00113-6}{doi:10.1016/s0272-6386(03)00113-6}. \url{https://pubmed.ncbi.nlm.nih.gov/12751050} (accessed 2026-11-15).
\noindent\hspace*{2em}\small\textit{Evidence:} ``Among the homeostatic processes controlling the endogenous L-carnitine pool in humans, the kidney has a vital role through extensive and adaptive tubular reabsorption.''
\noindent\hspace*{2em}\small\textit{Evidence:} ``Exogenous L-carnitine, administered intravenously, is approved for the treatment of secondary carnitine deficiency caused by long-term hemodialysis.''
\noindent [201] Fornasini G; Upton RN; Evans AM (2007). A pharmacokinetic model for L-carnitine in patients receiving haemodialysis.. \textit{Br J Clin Pharmacol} 64(3):335-45. \href{https://doi.org/10.1111/j.1365-2125.2007.02926.x}{doi:10.1111/j.1365-2125.2007.02926.x}. \url{https://pubmed.ncbi.nlm.nih.gov/17506778} (accessed 2026-11-15).
\noindent\hspace*{2em}\small\textit{Evidence:} ``Haemodialysis causes rapid reductions in plasma L-carnitine concentrations with each dialysis session.''
\noindent\hspace*{2em}\small\textit{Evidence:} ``the ongoing dialytic loss of L-carnitine may lead to a slow depletion of the compound, contributing to a possible secondary deficiency.''
\noindent [202] Görür S; Bağdatoğlu OT; Polat G (2005). Protective effect of L-carnitine on renal ischaemia-reperfusion injury in the rat.. \textit{Cell Biochem Funct} 23(3):151-5. \href{https://doi.org/10.1002/cbf.1159}{doi:10.1002/cbf.1159}. \url{https://pubmed.ncbi.nlm.nih.gov/15386530} (accessed 2026-11-15).
\noindent\hspace*{2em}\small\textit{Evidence:} ``It is shown that carnitine protects kidney tissue against ischaemia-reperfusion injury.''
\noindent\hspace*{2em}\small\textit{Evidence:} ``received L-carnitine hydrochloride (100 mg kg(-1)) intraperitoneally.''
\noindent [203] Ye J; Li J; Yu Y; Wei Q; Deng W; Yu L (2010). L-carnitine attenuates oxidant injury in HK-2 cells via ROS-mitochondria pathway.. \textit{Regul Pept} 161(1-3):58-66. \href{https://doi.org/10.1016/j.regpep.2009.12.024}{doi:10.1016/j.regpep.2009.12.024}. \url{https://pubmed.ncbi.nlm.nih.gov/20093144} (accessed 2026-11-15).
\noindent\hspace*{2em}\small\textit{Evidence:} ``These results suggested that L-carnitine could protect HK-2 cells from H(2)O(2)-induced injury through the inhibition of oxidative damage, mitochondria dysfunction and ultimately inhibition of cell apoptosis''
\noindent [204] Idrovo JP; Yang WL; Nicastro J; Coppa GF; Wang P (2012). Stimulation of carnitine palmitoyltransferase 1 improves renal function and attenuates tissue damage after ischemia/reperfusion.. \textit{J Surg Res} 177(1):157-64. \href{https://doi.org/10.1016/j.jss.2012.05.053}{doi:10.1016/j.jss.2012.05.053}. \url{https://pubmed.ncbi.nlm.nih.gov/22698429} (accessed 2026-11-15).
\noindent\hspace*{2em}\small\textit{Evidence:} ``Fatty acid β-oxidation is the major metabolic pathway for generating adenosine triphosphate (ATP) in the kidneys, which is governed by carnitine palmitoyltransferase 1 (CPT1).''
\noindent\hspace*{2em}\small\textit{Evidence:} ``Stimulation of CPT1 activity by C75 recovered ATP depletion, improved renal function, attenuated tissue injury, and inhibited proinflammatory cytokine production and neutrophil infiltration after renal I/R injury.''
\noindent [206] Erbaş M; Gündüz D; Sarman E; İlhan İ; Tepebaşı MY; Arlıoğlu M; Savran M (2026). L-carnitine mitigates urographin-induced nephrotoxicity via modulation of inflammatory, oxidative, and apoptotic pathways in rat kidney.. \textit{Mol Cell Biochem} 481(2):911-922. \href{https://doi.org/10.1007/s11010-025-05433-6}{doi:10.1007/s11010-025-05433-6}. \url{https://pubmed.ncbi.nlm.nih.gov/41261345} (accessed 2026-11-15).
\noindent\hspace*{2em}\small\textit{Evidence:} ``LC exerts potent nephroprotective effects against URO-induced renal injury by attenuating oxidative stress, suppressing pro-inflammatory and apoptotic signaling, and promoting anti-apoptotic gene expression.''
\noindent [207] Armaly Z; Artol S; Jabbour AR; Saffouri A; Habashi N; Abd Elkadir A; Ghattas N; Farah R; Kinaneh S; Nseir W (2019). Impact of pretreatment with carnitine and tadalafil on contrast-induced nephropathy in CKD patients.. \textit{Ren Fail} 41(1):976-986. \href{https://doi.org/10.1080/0886022X.2019.1669459}{doi:10.1080/0886022X.2019.1669459}. \url{https://pubmed.ncbi.nlm.nih.gov/31797710} (accessed 2026-11-15).
\noindent\hspace*{2em}\small\textit{Evidence:} ``In contrast, pretreatment with carnitine prevented the increase in urinary NGAL and reduced SCr below basal levels.''
\noindent [208] Mohammadi M; Hajhossein Talasaz A; Alidoosti M; Pour Hosseini HR; Gholami K; Jalali A; Aryannejad H (2017). Nephroprotective Effects of L-Carnitine against Contrast-Induced Nephropathy in Patients Undergoing Percutaneous Coronary Intervention: A Randomized Open-Labeled Clinical Trial.. \textit{J Tehran Heart Cent} 12(2):57-64. \href{https://pubmed.ncbi.nlm.nih.gov/28828020/}{PMID:28828020}. \url{https://pubmed.ncbi.nlm.nih.gov/28828020} (accessed 2026-11-15).
\noindent\hspace*{2em}\small\textit{Evidence:} ``Our results showed that oral L-carnitine was able to prevent an increase in NGAL following contrast medium administration in patients undergoing PCI.''
\noindent [209] Jones AE; Puskarich MA; Shapiro NI; Guirgis FW; Runyon M; Adams JY; Sherwin R; Arnold R; Roberts BW; Kurz MC; Wang HE; Kline JA; Courtney DM; Trzeciak S; Sterling SA; Nandi U; Patki D; Viele K (2018). Effect of Levocarnitine vs Placebo as an Adjunctive Treatment for Septic Shock: The Rapid Administration of Carnitine in Sepsis (RACE) Randomized Clinical Trial.. \textit{JAMA Netw Open} 1(8):e186076. \href{https://doi.org/10.1001/jamanetworkopen.2018.6076}{doi:10.1001/jamanetworkopen.2018.6076}. \url{https://pubmed.ncbi.nlm.nih.gov/30646314} (accessed 2026-11-15).
\noindent\hspace*{2em}\small\textit{Evidence:} ``In this dose-finding, phase 2 adaptive randomized trial, the most efficacious dose of levocarnitine (18 g) did not meaningfully reduce cumulative organ failure at 48 hours.''
\noindent [210] Puskarich MA; Kline JA; Krabill V; Claremont H; Jones AE (2014). Preliminary safety and efficacy of L-carnitine infusion for the treatment of vasopressor-dependent septic shock: a randomized control trial.. \textit{JPEN J Parenter Enteral Nutr} 38(6):736-43. \href{https://doi.org/10.1177/0148607113495414}{doi:10.1177/0148607113495414}. \url{https://pubmed.ncbi.nlm.nih.gov/23851424} (accessed 2026-11-15).
\noindent\hspace*{2em}\small\textit{Evidence:} ``Mortality was significantly lower at 28 days in the L-carnitine group (4/16 vs 9/15, P = .048), with a nonsignificant improved survival at 1 year (P = .06).''
\noindent [211] Guedes GV; Minicucci MF; Tanni SE (2022). The supplementation of L-carnitine in septic shock patients: Systematic review and meta-analysis.. \textit{Clinics (Sao Paulo)} 77:100124. \href{https://doi.org/10.1016/j.clinsp.2022.100124}{doi:10.1016/j.clinsp.2022.100124}. \url{https://pubmed.ncbi.nlm.nih.gov/36327640} (accessed 2026-11-15).
\noindent\hspace*{2em}\small\textit{Evidence:} ``There is low-quality evidence that the use of L-carnitine has no significant effect on reducing 28-day mortality in septic shock patients.''
\noindent [213] Jafari A; Khatami MR; Dashti-Khavidaki S; Lessan-Pezeshki M; Abdollahi A; Moghaddas A (2017). Protective Effects of L-Carnitine Against Delayed Graft Function in Kidney Transplant Recipients: A Pilot, Randomized, Double-Blinded, Placebo-Controlled Clinical Trial.. \textit{J Ren Nutr} 27(2):113-126. \href{https://doi.org/10.1053/j.jrn.2016.11.002}{doi:10.1053/j.jrn.2016.11.002}. \url{https://pubmed.ncbi.nlm.nih.gov/28065453} (accessed 2026-11-15).
\noindent\hspace*{2em}\small\textit{Evidence:} ``This study showed no protective effects of oral L-carnitine supplementation against DGF occurrence recipients; however, 3-month graft loss was lower in the L-carnitine supplemented group.''
\noindent [214] Gewin LS (2021). Sugar or Fat? Renal Tubular Metabolism Reviewed in Health and Disease.. \textit{Nutrients} 13(5). \href{https://doi.org/10.3390/nu13051580}{doi:10.3390/nu13051580}. \url{https://pubmed.ncbi.nlm.nih.gov/34065078} (accessed 2026-11-15).
\noindent\hspace*{2em}\small\textit{Evidence:} ``Thus, proximal tubules use fatty acid oxidation, which generates more adenosine triphosphate (ATP) than glucose metabolism, as its preferred metabolic pathway.''
\noindent [215] Portilla D (1999). Role of fatty acid beta-oxidation and calcium-independent phospholipase A2 in ischemic acute renal failure.. \textit{Curr Opin Nephrol Hypertens} 8(4):473-7. \href{https://doi.org/10.1097/00041552-199907000-00012}{doi:10.1097/00041552-199907000-00012}. \url{https://pubmed.ncbi.nlm.nih.gov/10491743} (accessed 2026-11-15).
\noindent\hspace*{2em}\small\textit{Evidence:} ``Altogether, lysophospholipids, long-chain acyl carnitines, and long-chain acyl coenzyme A inhibit proximal tubule Na+K(+)-ATPase.''
\noindent [216] Zhang M; Ma Y; Jin Y; Wang Y; Wu X (2025). Acute kidney injury and energy metabolism.. \textit{Clin Chim Acta} 570:120208. \href{https://doi.org/10.1016/j.cca.2025.120208}{doi:10.1016/j.cca.2025.120208}. \url{https://pubmed.ncbi.nlm.nih.gov/39986590} (accessed 2026-11-15).
\noindent\hspace*{2em}\small\textit{Evidence:} ``Special emphasis is placed on the role of fatty acid and carnitine metabolism in AKI, with the goal of providing a theoretical foundation for future investigations into AKI mechanisms and the identification of early diagnostic biomarkers.''
\noindent [217] Ito S; Taguchi K; Kodama G; Kubo S; Moriyama T; Yamashita Y; Yokota Y; Nakayama Y; Kaida Y; Shinohara M; Tashiro K; Ohta K; Yamagishi SI; Fukami K (2025). Involvement of impaired carnitine-induced fatty acid oxidation in experimental and human diabetic kidney disease.. \textit{JCI Insight} 10(13). \href{https://doi.org/10.1172/jci.insight.179362}{doi:10.1172/jci.insight.179362}. \url{https://pubmed.ncbi.nlm.nih.gov/40402578} (accessed 2026-11-15).
\noindent\hspace*{2em}\small\textit{Evidence:} ``In our single-center randomized control trial with patients undergoing peritoneal dialysis, l-carnitine supplementation preserved residual renal function and increased urine volume, the latter of which was correlated with improvement of tubular injury.''
\noindent [220] Li P; Wang Y; Luo J; Zeng Q; Wang M; Bai M; Zhou H; Wang J; Jiang H (2020). Downregulation of OCTN2 by cytokines plays an important role in the progression of inflammatory bowel disease.. \textit{Biochem Pharmacol} 178:114115. \href{https://doi.org/10.1016/j.bcp.2020.114115}{doi:10.1016/j.bcp.2020.114115}. \url{https://pubmed.ncbi.nlm.nih.gov/32579962} (accessed 2026-11-15).
\noindent\hspace*{2em}\small\textit{Evidence:} ``Mixed proinflammatory cytokines TNF-α, IL-1β and IFNγ downregulated the expression of OCTN2 and subsequently reduced the l-Car content through PPARγ/RXRα pathways in FHC cells.''
\noindent [221] Suzuki M; Kanaya M; Muramatsu S; Takahashi T (1976). Effects of carnitine administration, fasting, and exercise on urinary carnitine excretion in man.. \textit{J Nutr Sci Vitaminol (Tokyo)} 22(2):167-74. \href{https://pubmed.ncbi.nlm.nih.gov/956930/}{PMID:956930}. \url{https://pubmed.ncbi.nlm.nih.gov/956930} (accessed 2026-11-15).
\noindent\hspace*{2em}\small\textit{Evidence:} ``Urinary carnitine excretion was significantly increased during five-day fasting; a maximum level of 2.17 +/- 0.24 (mean and SD) mmoles/day was 4.6 times higher than the usual level of 0.47 +/- 0.10 mmoles/day (p less than 0.001).''
\noindent\hspace*{2em}\small\textit{Evidence:} ``The urinary excretion of carnitine increased slightly with heavy running exercise.''
\noindent [222] Arenas J; Ricoy JR; Encinas AR; Pola P; D'Iddio S; Zeviani M; Didonato S; Corsi M (1991). Carnitine in muscle, serum, and urine of nonprofessional athletes: effects of physical exercise, training, and L-carnitine administration.. \textit{Muscle Nerve} 14(7):598-604. \href{https://doi.org/10.1002/mus.880140703}{doi:10.1002/mus.880140703}. \url{https://pubmed.ncbi.nlm.nih.gov/1922166} (accessed 2026-11-15).
\noindent\hspace*{2em}\small\textit{Evidence:} ``These changes were likely related with the significantly increased urinary excretion of esterified and total carnitine which occurred after physical exercise.''
\noindent [223] Tosi I; Art T; Boemer F; Votion DM; Davis MS (2021). Acylcarnitine profile in Alaskan sled dogs during submaximal multiday exercise points out metabolic flexibility and liver role in energy metabolism.. \textit{PLoS One} 16(8):e0256009. \href{https://doi.org/10.1371/journal.pone.0256009}{doi:10.1371/journal.pone.0256009}. \url{https://pubmed.ncbi.nlm.nih.gov/34383825} (accessed 2026-11-15).
\noindent\hspace*{2em}\small\textit{Evidence:} ``Our results show chain-specific modification of AC profiles during the exercise challenge: LCACs maintained a steady increase throughout exercise, some SCACs increased during the last phase of exercise and acetylcarnitine (C2) initially increased before decreasing during the later phase of exercise.''
\noindent [225] Koeth RA; Wang Z; Levison BS; Buffa JA; Org E; Sheehy BT; Britt EB; Fu X; Wu Y; Li L; Smith JD; DiDonato JA; Chen J; Li H; Wu GD; Lewis JD; Warrier M; Brown JM; Krauss RM; Tang WH; Bushman FD; Lusis AJ; Hazen SL (2013). Intestinal microbiota metabolism of L-carnitine, a nutrient in red meat, promotes atherosclerosis. \textit{Nature medicine} 19(5):576-85. \href{https://doi.org/10.1038/nm.3145}{doi:10.1038/nm.3145}. \url{https://pmc.ncbi.nlm.nih.gov/articles/PMC3650111} (accessed 2026-11-15).
\noindent\hspace*{2em}\small\textit{Evidence:} ``Omnivorous subjects are shown to produce significantly more TMAO than vegans/vegetarians following ingestion of L-carnitine through a microbiota-dependent mechanism.''
\noindent\hspace*{2em}\small\textit{Evidence:} ``After a 12 hour overnight fast, subjects received a capsule of d3-carnitine (250 mg) alone, or in some cases (as in data for subject shown) also an 8 ounce steak (estimated 180 mg L-carnitine), whereupon serial plasma and 24h urine collection was obtained for TMA and TMAO analyses.''
\noindent\hspace*{2em}\small\textit{Evidence:} ``Chronic dietary L-carnitine supplementation in mice significantly altered cecal microbial composition, markedly enhanced synthesis of TMA/TMAO, and increased atherosclerosis, but not following suppression of intestinal microbiota.''
\noindent [227] Koeth RA; Lam-Galvez BR; Kirsop J; Wang Z; Levison BS; Gu X; Copeland MF; Bartlett D; Cody DB; Dai HJ; Culley MK; Li XS; Fu X; Wu Y; Li L; DiDonato JA; Tang WHW; Garcia-Garcia JC; Hazen SL (2019). l-Carnitine in omnivorous diets induces an atherogenic gut microbial pathway in humans. \textit{The Journal of clinical investigation} 129(1):373-387. \href{https://doi.org/10.1172/JCI94601}{doi:10.1172/JCI94601}. \url{https://pmc.ncbi.nlm.nih.gov/articles/PMC6307959} (accessed 2026-11-15).
\noindent\hspace*{2em}\small\textit{Evidence:} ``Moreover, a striking increase in d3-TMAO generation was observed in omnivores over vegans/vegetarians (>20-fold; \_P\_ = 0.001) following oral d3-l-carnitine ingestion, whereas fasting endogenous plasma l-carnitine and γBB levels were similar in vegans/vegetarians ( \_n\_ = 32) versus omnivores ( \_n\_ = 40).''
\noindent\hspace*{2em}\small\textit{Evidence:} ``Continuation of daily l-carnitine supplementation for at least an additional month resulted in no further increase in d3-l-carnitine→d3-TMAO transformation in omnivores, but continued to increase d3-TMAO generation in vegans/vegetarians following oral d3-l-carnitine challenge ( Figure 4A).''
\noindent\hspace*{2em}\small\textit{Evidence:} ``Gut microbiota γBB→TMA/TMAO transformation is induced by omnivorous dietary patterns and chronic l-carnitine exposure.''
\noindent [228] Wu WK; Chen CC; Liu PY; Panyod S; Liao BY; Chen PC; Kao HL; Kuo HC; Kuo CH; Chiu THT; Chen RA; Chuang HL; Huang YT; Zou HB; Hsu CC; Chang TY; Lin CL; Ho CT; Yu HT; Sheen LY; Wu MS (2019). Identification of TMAO-producer phenotype and host-diet-gut dysbiosis by carnitine challenge test in human and germ-free mice. \textit{Gut} 68(8):1439-1449. \href{https://doi.org/10.1136/gutjnl-2018-317155}{doi:10.1136/gutjnl-2018-317155}. \url{https://pmc.ncbi.nlm.nih.gov/articles/PMC6691853} (accessed 2026-11-15).
\noindent\hspace*{2em}\small\textit{Evidence:} ``The omnivores exhibited a 10-fold higher OR to be high TMAO producer than vegetarians.''
\noindent [229] Wu WK; Panyod S; Liu PY; Chen CC; Kao HL; Chuang HL; Chen YH; Zou HB; Kuo HC; Kuo CH; Liao BY; Chiu THT; Chung CH; Lin AY; Lee YC; Tang SL; Wang JT; Wu YW; Hsu CC; Sheen LY; Orekhov AN; Wu MS (2020). Characterization of TMAO productivity from carnitine challenge facilitates personalized nutrition and microbiome signatures discovery. \textit{Microbiome} 8(1):162. \href{https://doi.org/10.1186/s40168-020-00912-y}{doi:10.1186/s40168-020-00912-y}. \url{https://pmc.ncbi.nlm.nih.gov/articles/PMC7676756} (accessed 2026-11-15).
\noindent\hspace*{2em}\small\textit{Evidence:} ``Production of atherogenic trimethylamine N-oxide (TMAO) from carnitine is a gut microbiota-directed pathway and varies widely among individuals.''
\noindent\hspace*{2em}\small\textit{Evidence:} ``A microbiome-based random forest classifier was therefore constructed to predict the TMAO producer phenotype (AUROC = 0.81) which was then validated with an external cohort (AUROC = 0.80).''
\noindent [230] Missailidis C; Hallqvist J; Qureshi AR; Barany P; Heimburger O; Lindholm B; Stenvinkel P; Bergman P (2016). Serum Trimethylamine-N-Oxide Is Strongly Related to Renal Function and Predicts Outcome in Chronic Kidney Disease. \textit{PloS one} 11(1):e0141738. \href{https://doi.org/10.1371/journal.pone.0141738}{doi:10.1371/journal.pone.0141738}. \url{https://pubmed.ncbi.nlm.nih.gov/26751065} (accessed 2026-11-15).
\noindent\hspace*{2em}\small\textit{Evidence:} ``GFR was the dominant variable affecting TMAO (β = -0.41; p<0.001), choline (β = -0.38; p<0.001), and betaine (β = 0.45; p<0.001) levels.''
\noindent\hspace*{2em}\small\textit{Evidence:} ``Elevated TMAO levels are strongly associated with degree of renal function in CKD and normalize after renal transplantation.''
\noindent\hspace*{2em}\small\textit{Evidence:} ``A longitudinal study of 74 CKD 5 patients starting renal replacement therapy demonstrated that whereas dialysis treatment did not affect TMAO, Rtx reduced levels of TMAO to that of controls (p<0.001).''
\noindent [231] Stubbs JR; House JA; Ocque AJ; Zhang S; Johnson C; Kimber C; Schmidt K; Gupta A; Wetmore JB; Nolin TD; Spertus JA; Yu AS (2016). Serum Trimethylamine-N-Oxide is Elevated in CKD and Correlates with Coronary Atherosclerosis Burden. \textit{Journal of the American Society of Nephrology : JASN} 27(1):305-13. \href{https://doi.org/10.1681/ASN.2014111063}{doi:10.1681/ASN.2014111063}. \url{https://pubmed.ncbi.nlm.nih.gov/26229137} (accessed 2026-11-15).
\noindent\hspace*{2em}\small\textit{Evidence:} ``TMAO concentrations were markedly higher in patients receiving dialysis (median \textbackslash\{\}[interquartile range\textbackslash\{\}], 94.4 μM \textbackslash\{\}[54.8-133.0 μM\textbackslash\{\}] for dialysis-dependent patients versus 3.3 μM \textbackslash\{\}[3.1-6.0 μM\textbackslash\{\}] for healthy controls; P<0.001); whereas renal transplantation resulted in substantial reductions in TMAO concentrations (median \textbackslash\{\}[min-max\textbackslash\{\}] 71.2 μM \textbackslash\{\}[29.2-189.7 μM\textbackslash\{\}] pretransplant versus 11.4 μM \textbackslash\{\}[8.9-20.2 μM\textbackslash\{\}] post-transplant; P=0.03).''
\noindent [232] Tang WH; Wang Z; Kennedy DJ; Wu Y; Buffa JA; Agatisa-Boyle B; Li XS; Levison BS; Hazen SL (2015). Gut microbiota-dependent trimethylamine N-oxide (TMAO) pathway contributes to both development of renal insufficiency and mortality risk in chronic kidney disease. \textit{Circulation research} 116(3):448-55. \href{https://doi.org/10.1161/CIRCRESAHA.116.305360}{doi:10.1161/CIRCRESAHA.116.305360}. \url{https://pmc.ncbi.nlm.nih.gov/articles/PMC4312512} (accessed 2026-11-15).
\noindent\hspace*{2em}\small\textit{Evidence:} ``Within CKD subjects, higher (4th vs. 1st quartile) plasma TMAO level was associated with a 2.8-fold increased mortality risk. Following adjustments for traditional risk factors, hsCRP and eGFR, elevated TMAO levels remained predictive of 5-year mortality risk (HR 1.93 \textbackslash\{\}[95\%CI 1.13–3.29\textbackslash\{\}], p<0.05).''
\noindent\hspace*{2em}\small\textit{Evidence:} ``In animal models, elevated dietary choline or TMAO directly led to progressive renal tubulointerstitial fibrosis and dysfunction.''
\noindent [233] Wang M; Tang WHW; Li XS; de Oliveira Otto MC; Lee Y; Lemaitre RN; Fretts A; Nemet I; Sotoodehnia N; Sitlani CM; Budoff M; DiDonato JA; Wang Z; Bansal N; Shlipak MG; Psaty BM; Siscovick DS; Sarnak MJ; Mozaffarian D; Hazen SL (2024). The Gut Microbial Metabolite Trimethylamine N -oxide, Incident CKD, and Kidney Function Decline. \textit{Journal of the American Society of Nephrology : JASN} 35(6):749-760. \href{https://doi.org/10.1681/ASN.0000000000000344}{doi:10.1681/ASN.0000000000000344}. \url{https://pubmed.ncbi.nlm.nih.gov/38593157} (accessed 2026-11-15).
\noindent\hspace*{2em}\small\textit{Evidence:} ``Higher TMAO levels were associated with higher risk of incident CKD (second to fifth versus first quintile hazard ratio \textbackslash\{\}[95\% confidence interval\textbackslash\{\}]=1.65 \textbackslash\{\}[1.22 to 2.23\textbackslash\{\}], 1.68 \textbackslash\{\}[1.26 to 2.25\textbackslash\{\}], 2.28 \textbackslash\{\}[1.72 to 3.02\textbackslash\{\}], and 2.24 \textbackslash\{\}[1.68 to 2.98\textbackslash\{\}], respectively) and greater annualized eGFR decline (second to fifth versus first quintile annualized eGFR change=−0.21 \textbackslash\{\}[−0.32 to −0.09\textbackslash\{\}], −0.17 \textbackslash\{\}[−0.29 to −0.05\textbackslash\{\}], −0.35 \textbackslash\{\}[−0.47 to −0.22\textbackslash\{\}], and −0.43 \textbackslash\{\}[−0.56 to −0.30\textbackslash\{\}] ml/min per 1.73 m2, respectively) with monotonic dose–response relationships.''
\noindent [234] Stubbs JR; Stedman MR; Liu S; Long J; Franchetti Y; West RE 3rd; Prokopienko AJ; Mahnken JD; Chertow GM; Nolin TD (2019). Trimethylamine N-Oxide and Cardiovascular Outcomes in Patients with ESKD Receiving Maintenance Hemodialysis. \textit{Clinical journal of the American Society of Nephrology : CJASN} 14(2):261-267. \href{https://doi.org/10.2215/CJN.06190518}{doi:10.2215/CJN.06190518}. \url{https://pubmed.ncbi.nlm.nih.gov/30665924} (accessed 2026-11-15).
\noindent\hspace*{2em}\small\textit{Evidence:} ``We observed a wide distribution of serum TMAO in our cohort, with approximately 80\% of participants exhibiting TMAO concentrations ≥56 \_µ\_ M and a maximum TMAO concentration of 1103.1 \_µ\_ M.''
\noindent\hspace*{2em}\small\textit{Evidence:} ``We found no evidence linking TMAO to adverse clinical outcomes in patients receiving maintenance hemodialysis with moderate to severe secondary hyperparathyroidism.''
\noindent [235] Kapetanaki S; Kumawat AK; Persson K; Demirel I (2021). The Fibrotic Effects of TMAO on Human Renal Fibroblasts Is Mediated by NLRP3, Caspase-1 and the PERK/Akt/mTOR Pathway. \textit{International journal of molecular sciences} 22(21):. \href{https://doi.org/10.3390/ijms222111864}{doi:10.3390/ijms222111864}. \url{https://pubmed.ncbi.nlm.nih.gov/34769294} (accessed 2026-11-15).
\noindent\hspace*{2em}\small\textit{Evidence:} ``We found that TMAO promoted renal fibroblast activation and fibroblast proliferation via the PERK/Akt/mTOR pathway, NLRP3, and caspase-1 signaling.''
\noindent\hspace*{2em}\small\textit{Evidence:} ``A high TMAO-containing diet has been found to cause tubulointerstitial renal fibrosis in mice.''
\noindent [238] Tang WH; Wang Z; Levison BS; Koeth RA; Britt EB; Fu X; Wu Y; Hazen SL (2013). Intestinal microbial metabolism of phosphatidylcholine and cardiovascular risk. \textit{The New England journal of medicine} 368(17):1575-84. \href{https://doi.org/10.1056/NEJMoa1109400}{doi:10.1056/NEJMoa1109400}. \url{https://pubmed.ncbi.nlm.nih.gov/23614584} (accessed 2026-11-15).
\noindent\hspace*{2em}\small\textit{Evidence:} ``Increased plasma levels of TMAO were associated with an increased risk of a major adverse cardiovascular event (hazard ratio for highest vs. lowest TMAO quartile, 2.54; 95\% confidence interval, 1.96 to 3.28; P<0.001).''
\noindent\hspace*{2em}\small\textit{Evidence:} ``Plasma levels of TMAO were markedly suppressed after the administration of antibiotics and then reappeared after withdrawal of antibiotics.''
\noindent [239] Schiattarella GG; Sannino A; Toscano E; Giugliano G; Gargiulo G; Franzone A; Trimarco B; Esposito G; Perrino C (2017). Gut microbe-generated metabolite trimethylamine-N-oxide as cardiovascular risk biomarker: a systematic review and dose-response meta-analysis. \textit{European heart journal} 38(39):2948-2956. \href{https://doi.org/10.1093/eurheartj/ehx342}{doi:10.1093/eurheartj/ehx342}. \url{https://pubmed.ncbi.nlm.nih.gov/29020409} (accessed 2026-11-15).
\noindent\hspace*{2em}\small\textit{Evidence:} ``High TMAO plasma levels were associated with increased incidence of all-cause mortality \textbackslash\{\}[14 studies for 16 cohorts enrolling 15 662 subjects, hazard ratio (HR): 1.91; 95\% confidence interval (CI): 1.40-2.61, P < 0.0001, I2 = 94\%\textbackslash\{\}] and MACCE (5 studies for 6 cohorts enrolling 13 944 subjects, HR: 1.67, 95\% CI: 1.33-2.11, P < 0.00001, I2 = 46\%,). Dose-response meta-analysis revealed that the relative risk (RR) for all-cause mortality increased by 7.6\% per each 10 μmol/L increment of TMAO \textbackslash\{\}[summary RR: 1.07, 95\% CI (1.04-1.11), P < 0.0001; based on seven studies\textbackslash\{\}].''
\noindent [240] Budoff MJ; de Oliveira Otto MC; Li XS; Lee Y; Wang M; Lai HTM; Lemaitre RN; Pratt A; Tang WHW; Psaty BM; Siscovick DS; Hazen SL; Mozaffarian D (2025). Trimethylamine-N-oxide (TMAO) and risk of incident cardiovascular events in the multi ethnic study of Atherosclerosis. \textit{Scientific reports} 15(1):23362. \href{https://doi.org/10.1038/s41598-025-05903-3}{doi:10.1038/s41598-025-05903-3}. \url{https://pubmed.ncbi.nlm.nih.gov/40603925} (accessed 2026-11-15).
\noindent\hspace*{2em}\small\textit{Evidence:} ``After multivariate adjustment, TMAO associated with higher risk of ASCVD in a dose-dependent fashion, with hazard ratios across quintiles of 1.02, 1.17, 1.23, and 1.33 (95\% CI 1.02, 1.74), respectively, compared to the lowest quintile (P-trend = 0.01).''
\noindent\hspace*{2em}\small\textit{Evidence:} ``Risk appeared possibly larger among Hispanic and Chinese adults; and among individuals with lower baseline renal function; although these interactions did not achieve statistical significance.''
\noindent [241] Jia J; Dou P; Gao M; Kong X; Li C; Liu Z; Huang T (2019). Assessment of Causal Direction Between Gut Microbiota-Dependent Metabolites and Cardiometabolic Health: A Bidirectional Mendelian Randomization Analysis. \textit{Diabetes} 68(9):1747-1755. \href{https://doi.org/10.2337/db19-0153}{doi:10.2337/db19-0153}. \url{https://pubmed.ncbi.nlm.nih.gov/31167879} (accessed 2026-11-15).
\noindent\hspace*{2em}\small\textit{Evidence:} ``Genetically predicted higher TMAO and carnitine were not associated with higher odds of T2DM, AF, CAD, MI, stroke, and CKD after Bonferroni correction ( \_P\_ ≤ 0.0005).''
\noindent\hspace*{2em}\small\textit{Evidence:} ``We further found that T2DM (0.130 ± 0.036, \_P\_ < 0.0001) and CKD (0.483 ± 0.168, \_P\_ = 0.004) were causally associated with higher TMAO levels.''
\noindent\hspace*{2em}\small\textit{Evidence:} ``Our Mendelian randomization findings support that T2DM and kidney disease increase TMAO levels and that observational evidence for cardiovascular diseases may be due to confounding or reverse causality.''
\noindent [242] Collins HL; Drazul-Schrader D; Sulpizio AC; Koster PD; Williamson Y; Adelman SJ; Owen K; Sanli T; Bellamine A (2016). L-Carnitine intake and high trimethylamine N-oxide plasma levels correlate with low aortic lesions in ApoE(-/-) transgenic mice expressing CETP. \textit{Atherosclerosis} 244:29-37. \href{https://doi.org/10.1016/j.atherosclerosis.2015.10.108}{doi:10.1016/j.atherosclerosis.2015.10.108}. \url{https://pubmed.ncbi.nlm.nih.gov/26584136} (accessed 2026-11-15).
\noindent\hspace*{2em}\small\textit{Evidence:} ``Surprisingly, and independently from treatment group, TMAO levels inversely correlated with aortic lesion size in both aortic root and thoracic aorta. High TMAO levels were found to significantly correlate with smaller aortic lesion area. Plasma lipid and lipoprotein levels did not change with treatment nor with TMAO levels, suggesting that the observed effects on lesion area were independent from lipid changes. Conclusion: These findings suggest that TMAO slows aortic lesion formation in this mouse model and may have a protective effect against atherosclerosis development in humans.''
\noindent [243] Johri AM; Hetu MF; Heyland DK; Herr JE; Korol J; Froese S; Norman PA; Day AG; Matangi MF; Michos ED; LaHaye SA; Saunders FW; Spence JD (2022). Progression of atherosclerosis with carnitine supplementation: a randomized controlled trial in the metabolic syndrome. \textit{Nutrition \& metabolism} 19(1):26. \href{https://doi.org/10.1186/s12986-022-00661-9}{doi:10.1186/s12986-022-00661-9}. \url{https://pubmed.ncbi.nlm.nih.gov/35366920} (accessed 2026-11-15).
\noindent\hspace*{2em}\small\textit{Evidence:} ``In 157 participants (L-C N = 76, placebo N = 81), no difference in TPV change between arms was found. The L-C group had a greater increase in carotid atherosclerotic stenosis of 9.3\% (p = 0.02) than the placebo group.''
\noindent [244] Wang Z; Tang WHW; O'Connell T; Garcia E; Jeyarajah EJ; Li XS; Jia X; Weeks TL; Hazen SL (2022). Circulating trimethylamine N-oxide levels following fish or seafood consumption. \textit{European journal of nutrition} 61(5):2357-2364. \href{https://doi.org/10.1007/s00394-022-02803-4}{doi:10.1007/s00394-022-02803-4}. \url{https://pmc.ncbi.nlm.nih.gov/articles/PMC9283263} (accessed 2026-11-15).
\noindent\hspace*{2em}\small\textit{Evidence:} ``Consumption of fish sticks, cod, and to lesser extent, salmon, lead to significant increases in circulating TMAO levels. Within one day, circulating TMAO concentrations in all participants returned to baseline levels.''
\noindent\hspace*{2em}\small\textit{Evidence:} ``The likelihood of having an elevated (>6.4 μM) TMAO level the morning after eating either fish or red meat was 45\% or 20\%, respectively. Among participants that ate neither fish nor red meat the previous night, the likelihood of elevated TMAO level at blood draw was 7.4\%.''
\noindent\hspace*{2em}\small\textit{Evidence:} ``We conclude that some fish and seafood contain high levels of TMAO, and may induce a transient elevation in TMAO levels in some individuals.''
\noindent [245] Costabile G; Vetrani C; Bozzetto L; Giacco R; Bresciani L; Del Rio D; Vitale M; Della Pepa G; Brighenti F; Riccardi G; Rivellese AA; Annuzzi G (2021). Plasma TMAO increase after healthy diets: results from 2 randomized controlled trials with dietary fish, polyphenols, and whole-grain cereals. \textit{The American journal of clinical nutrition} 114(4):1342-1350. \href{https://doi.org/10.1093/ajcn/nqab188}{doi:10.1093/ajcn/nqab188}. \url{https://pubmed.ncbi.nlm.nih.gov/34091663} (accessed 2026-11-15).
\noindent\hspace*{2em}\small\textit{Evidence:} ``Diets rich in LCn3 of marine origin or WGCs significantly increased plasma TMAO concentration. These changes mirrored the direct associations between TMAO concentrations and intakes of fish and WGCs, suggesting that TMAO reflects intakes of these healthy foods and, therefore, it is not a universally valid biomarker of cardiometabolic risk independent of the background diet.''
\noindent\hspace*{2em}\small\textit{Evidence:} ``Among food groups, TMAO directly correlated with the intake of fish, vegetables, and whole-grain products, but not meat, processed meat, and dairy products.''
\noindent [246] Davani-Davari D; Karimzadeh I; Sagheb MM; Khalili H (2019). The Renal Safety of L-Carnitine, L-Arginine, and Glutamine in Athletes and Bodybuilders. \textit{Journal of renal nutrition : the official journal of the Council on Renal Nutrition of the National Kidney Foundation} 29(3):221-234. \href{https://doi.org/10.1053/j.jrn.2018.08.014}{doi:10.1053/j.jrn.2018.08.014}. \url{https://pubmed.ncbi.nlm.nih.gov/30341034} (accessed 2026-11-15).
\noindent\hspace*{2em}\small\textit{Evidence:} ``The risks and benefits of L-carnitine on the athletes' and bodybuilders' kidney have not been evaluated yet. However, L-carnitine up to 6000 mg/day is generally considered to be a safe supplement at least in healthy adults.''
\noindent [247] Olek RA; Radak Z (2026). Trimethylamine N-oxide in exercise physiology: a gut microbiota-derived signal linking metabolic stress, redox balance and cardiometabolic health. \textit{European journal of applied physiology} 126(8):4229-4242. \href{https://doi.org/10.1007/s00421-026-06297-4}{doi:10.1007/s00421-026-06297-4}. \url{https://pubmed.ncbi.nlm.nih.gov/42287393} (accessed 2026-11-15).
\noindent\hspace*{2em}\small\textit{Evidence:} ``Across human studies, TMAO levels do not consistently correlate with aerobic fitness, habitual physical activity, or training status. Exercise-only interventions rarely produce significant changes in circulating TMAO concentrations, whereas dietary manipulation and precursor supplementation exert substantially stronger effects.''
\noindent [248] Hathcock JN; Shao A (2006). Risk assessment for carnitine. \textit{Regulatory toxicology and pharmacology : RTP} 46(1):23-8. \href{https://doi.org/10.1016/j.yrtph.2006.06.007}{doi:10.1016/j.yrtph.2006.06.007}. \url{https://pubmed.ncbi.nlm.nih.gov/16901595} (accessed 2026-11-15).
\noindent\hspace*{2em}\small\textit{Evidence:} ``The OSL risk assessment method indicates that the evidence of safety is strong at intakes up to 2000mg/day l-carnitine equivalents for chronic supplementation, and this level is identified as the OSL.''
\noindent\hspace*{2em}\small\textit{Evidence:} ``Although much higher levels have been tested without adverse effects and may be safe, the data for intakes above 2000mg/day are not sufficient for a confident conclusion of long-term safety.''
\noindent [249] Papandreou C; More M; Bellamine A (2020). Trimethylamine N-Oxide in Relation to Cardiometabolic Health-Cause or Effect?. \textit{Nutrients} 12(5):. \href{https://doi.org/10.3390/nu12051330}{doi:10.3390/nu12051330}. \url{https://pubmed.ncbi.nlm.nih.gov/32392758} (accessed 2026-11-15).
\noindent\hspace*{2em}\small\textit{Evidence:} ``Several animal studies reported neutral or even beneficial effects of TMAO or its precursors in cardiovascular disease model systems, supporting the clinically proven beneficial effects of its precursor, L-carnitine, or a sea-food rich diet (naturally containing TMAO) on cardiometabolic health.''
\noindent\hspace*{2em}\small\textit{Evidence:} ``Fish intake provides a direct significant source of TMAO.''
\noindent [253] Bachmann HU; Hoffmann A (2004). Interaction of food supplement L-carnitine with oral anticoagulant acenocoumarol. \textit{Swiss Med Wkly} 134(25-26):385. \href{https://doi.org/10.4414/smw.2004.10631}{doi:10.4414/smw.2004.10631}. \url{https://smw.ch/index.php/smw/article/download/389/386/766} (accessed 2026-11-15).
\noindent\hspace*{2em}\small\textit{Evidence:} ``An unexpected rise of INR to 4.65 which was observed again despite a dose correction was the reason for a detailed questioning on food intake.''
\noindent\hspace*{2em}\small\textit{Evidence:} ``The patient had been taking 1000 mg/d of carnitine for a period of ten weeks.''
\noindent\hspace*{2em}\small\textit{Evidence:} ``after the product was omitted the INR values returned to a stable range (2.4–2.8) for over 4 months.''
\noindent [255] Martinez E; Domingo P; Roca-Cusachs A (1993). Potentiation of acenocoumarol action by L-carnitine.. \textit{J Intern Med} 233(1):94. \href{https://pubmed.ncbi.nlm.nih.gov/8429297/}{PMID:8429297}. \url{https://pubmed.ncbi.nlm.nih.gov/8429297} (accessed 2026-11-15).
\noindent\hspace*{2em}\small\textit{Evidence:} ``Potentiation of acenocoumarol action by L-carnitine.''
\noindent [256] Veronese N; Stubbs B; Solmi M; Ajnakina O; Carvalho AF; Maggi S (2018). Acetyl-L-Carnitine Supplementation and the Treatment of Depressive Symptoms: A Systematic Review and Meta-Analysis.. \textit{Psychosom Med} 80(2):154-159. \href{https://doi.org/10.1097/PSY.0000000000000537}{doi:10.1097/PSY.0000000000000537}. \url{https://pubmed.ncbi.nlm.nih.gov/29076953} (accessed 2026-11-15).
\noindent\hspace*{2em}\small\textit{Evidence:} ``ALC supplementation significantly decreases depressive symptoms compared with placebo/no intervention, while offering a comparable effect with that of established antidepressant agents with fewer adverse effects.''
\noindent\hspace*{2em}\small\textit{Evidence:} ``Twelve RCTs (11 of which were ALC monotherapy) with a total of 791 participants''
\noindent\hspace*{2em}\small\textit{Evidence:} ``the incidence of adverse effects was significantly lower in the ALC group than in the antidepressant group.''
\noindent [257] Bazzato G; Mezzina C; Ciman M; Guarnieri G (1979). Myasthenia-like syndrome associated with carnitine in patients on long-term haemodialysis.. \textit{Lancet} 1(8124):1041-2. \href{https://doi.org/10.1016/s0140-6736(79)92804-6}{doi:10.1016/s0140-6736(79)92804-6}. \url{https://pubmed.ncbi.nlm.nih.gov/86768} (accessed 2026-11-15).
\noindent\hspace*{2em}\small\textit{Evidence:} ``Myasthenia-like syndrome associated with carnitine in patients on long-term haemodialysis.''
\noindent [258] Benvenga S; Ruggeri RM; Russo A; Lapa D; Campenni A; Trimarchi F (2001). Usefulness of L-carnitine, a naturally occurring peripheral antagonist of thyroid hormone action, in iatrogenic hyperthyroidism: a randomized, double-blind, placebo-controlled clinical trial.. \textit{J Clin Endocrinol Metab} 86(8):3579-94. \href{https://doi.org/10.1210/jcem.86.8.7747}{doi:10.1210/jcem.86.8.7747}. \url{https://pubmed.ncbi.nlm.nih.gov/11502782} (accessed 2026-11-15).
\noindent\hspace*{2em}\small\textit{Evidence:} ``In conclusion, L-carnitine is effective in both reversing and preventing symptoms of hyperthyroidism and has a beneficial effect on bone mineralization.''
\noindent\hspace*{2em}\small\textit{Evidence:} ``we assessed whether 2 or 4 g/d oral L-carnitine were able to both reverse and prevent/minimize nine hyperthyroidism- related symptoms''
\noindent\hspace*{2em}\small\textit{Evidence:} ``Because hyperthyroidism depletes the body deposits of carnitine and since carnitine has no toxicity, teratogenicity, contraindications and interactions with drugs, carnitine can be of clinical use.''
\noindent [259] Benvenga S; Amato A; Calvani M; Trimarchi F (2004). Effects of carnitine on thyroid hormone action.. \textit{Ann N Y Acad Sci} 1033:158-67. \href{https://doi.org/10.1196/annals.1320.015}{doi:10.1196/annals.1320.015}. \url{https://pubmed.ncbi.nlm.nih.gov/15591013} (accessed 2026-11-15).
\noindent\hspace*{2em}\small\textit{Evidence:} ``In particular, L-carnitine inhibits both triiodothyronine (T3) and thyroxine (T4) entry into the cell nuclei.''
\noindent\hspace*{2em}\small\textit{Evidence:} ``we showed that 2 and 4 grams per day of oral L-carnitine are capable of reversing hyperthyroid symptoms (and biochemical changes in the hyperthyroid direction) as well as preventing (or minimizing) the appearance of hyperthyroid symptoms''
\noindent [260] Troutman AD; Srinivasan S; Metzger CE; Fallen PB; Chen N; O'Neill KD; Allen MR; Biruete A; Moe SM; Avin KG (2024). Musculoskeletal Health Worsened from Carnitine Supplementation and Not Impacted by a Novel Individualized Treadmill Training Protocol.. \textit{Am J Nephrol} 55(3):369-379. \href{https://doi.org/10.1159/000537827}{doi:10.1159/000537827}. \url{https://pubmed.ncbi.nlm.nih.gov/38377965} (accessed 2026-11-15).
\noindent\hspace*{2em}\small\textit{Evidence:} ``Carnitine supplementation worsened CKD progression, mineral metabolism biochemistries, and cortical porosity and did not have an impact on physical function.''
\noindent\hspace*{2em}\small\textit{Evidence:} ``Future studies should seek to better understand carnitine doses in conditions of compromised renal function to prevent toxicity which may result from elevated carnitine levels''
\noindent [261] Clair F; Caillat S; Soufir JC; Lafforgue B; Drüeke T; Saïd G (1984). [Myasthenic syndrome induced by D,L-carnitine in a chronic hemodialysis patient].. \textit{Presse Med} 13(18):1154-5. \href{https://pubmed.ncbi.nlm.nih.gov/6232557/}{PMID:6232557}. \url{https://pubmed.ncbi.nlm.nih.gov/6232557} (accessed 2026-11-15).
\noindent\hspace*{2em}\small\textit{Evidence:} ``[Myasthenic syndrome induced by D,L-carnitine in a chronic hemodialysis patient].''
\noindent [262] De Grandis D; Mezzina C; Fiaschi A; Pinelli P; Bazzato G; Morachiello M (1980). Myasthenia due to carnitine treatment.. \textit{J Neurol Sci} 46(3):365-71. \href{https://doi.org/10.1016/0022-510x(80)90061-1}{doi:10.1016/0022-510x(80)90061-1}. \url{https://pubmed.ncbi.nlm.nih.gov/7381520} (accessed 2026-11-15).
\noindent\hspace*{2em}\small\textit{Evidence:} ``Severe weakness of lower cranial and girdle muscles has been found in 4 subjects among 20 dialysis patients treated with carnitine in order to correct their high plasma triglyceride level.''
\noindent\hspace*{2em}\small\textit{Evidence:} ``These findings suggest a hemicholinium activity of carnitine at synaptic level.''
\noindent [264] Jacobs PL; Goldstein ER; Blackburn W; Orem I; Hughes JJ (2009). Glycine propionyl-L-carnitine produces enhanced anaerobic work capacity with reduced lactate accumulation in resistance trained males.. \textit{J Int Soc Sports Nutr} 6:9. \href{https://doi.org/10.1186/1550-2783-6-9}{doi:10.1186/1550-2783-6-9}. \url{https://pubmed.ncbi.nlm.nih.gov/19341458} (accessed 2026-11-15).
\noindent\hspace*{2em}\small\textit{Evidence:} ``Testing was performed 90 minutes following oral ingestion of either 4.5 grams GPLC or 4.5 grams cellulose (PL), in randomized order.''
\noindent\hspace*{2em}\small\textit{Evidence:} ``These findings indicate that short-term oral supplementation of GPLC can enhance peak power production in resistance trained males with significantly less LAC accumulation.''
\noindent [265] Yahyapoor F; Keshani M; Sedaghat A; Feizi A; Clark CCT; Bagherniya M; Safarian M; Badpeyma M; Norouzy A (2023). The effects of adjunctive treatment with L-carnitine on monitoring laboratory variables in ICU patients: a double-blinded randomized controlled clinical trial.. \textit{Trials} 24(1):3. \href{https://doi.org/10.1186/s13063-022-07010-4}{doi:10.1186/s13063-022-07010-4}. \url{https://pubmed.ncbi.nlm.nih.gov/36597167} (accessed 2026-11-15).
\noindent\hspace*{2em}\small\textit{Evidence:} ``Alanine transaminase (ALT) (0.022), lactate (<0.001), creatinine (Cr) (0.005), and international normalized ratio (INR) (0.049) decreased meaningfully in the intervention vs. control group.''
\noindent\hspace*{2em}\small\textit{Evidence:} ``Participants were 54 patients, aged > 18 years, with multiple conditions, randomly assigned to receive 3 g L-carnitine per day or placebo, along with enteral feeding, for 1 week.''
\noindent [266] Krims-Davis K; Ozola M; Razzivina V; Gukalova B; Konrade I; Dambrova M; Liepinsh E (2025). Low Bioavailability and High TMAO Production: Novel Insights Into Acetylcarnitine and Carnitine Metabolism.. \textit{Mol Nutr Food Res} 69(24):e70316. \href{https://doi.org/10.1002/mnfr.70316}{doi:10.1002/mnfr.70316}. \url{https://pubmed.ncbi.nlm.nih.gov/41243468} (accessed 2026-11-15).
\noindent\hspace*{2em}\small\textit{Evidence:} ``Healthy volunteers received either carnitine or an acetylcarnitine supplement (0.5 or 1.5 g).''
\noindent\hspace*{2em}\small\textit{Evidence:} ``Acetylcarnitine has significantly lower bioavailability than carnitine.''
\noindent\hspace*{2em}\small\textit{Evidence:} ``Approximately 90\% of both supplements were metabolized to TMAO, reaching 50 µM in plasma-levels previously found to be associated with adverse health outcomes.''
\noindent [267] Kraemer WJ; Volek JS; French DN; Rubin MR; Sharman MJ; Gómez AL; Ratamess NA; Newton RU; Jemiolo B; Craig BW; Häkkinen K (2003). The effects of L-carnitine L-tartrate supplementation on hormonal responses to resistance exercise and recovery.. \textit{J Strength Cond Res} 17(3):455-62. \href{https://doi.org/10.1519/1533-4287(2003)017<0455:teolls>2.0.co;2}{doi:10.1519/1533-4287(2003)017<0455:teolls>2.0.co;2}. \url{https://pubmed.ncbi.nlm.nih.gov/12930169} (accessed 2026-11-15).
\noindent\hspace*{2em}\small\textit{Evidence:} ``after 3 weeks of supplementation (2 g LCLT per day), fasting morning blood samples were obtained on six consecutive days''
\noindent\hspace*{2em}\small\textit{Evidence:} ``Over the recovery period, LCLT reduced the amount of exercise-induced muscle tissue damage, which was assessed via magnetic resonance imaging scans of the thigh.''
\noindent [269] Bain MA; Milne RW; Evans AM (2006). Disposition and metabolite kinetics of oral L-carnitine in humans.. \textit{J Clin Pharmacol} 46(10):1163-70. \href{https://doi.org/10.1177/0091270006292851}{doi:10.1177/0091270006292851}. \url{https://pubmed.ncbi.nlm.nih.gov/16988205} (accessed 2026-11-15).
\noindent\hspace*{2em}\small\textit{Evidence:} ``The pharmacokinetics of L-carnitine and its metabolites were investigated in 7 healthy subjects following the oral administration of 0, 0.5, 1, and 2 g 3 times a day for 7 days.''
\noindent\hspace*{2em}\small\textit{Evidence:} ``Mean plasma concentrations of L-carnitine across an 8-hour dose interval increased significantly (P < .05) from a baseline of 54.2 +/- 9.3 microM to 80.5 +/- 12.5 microM following the 0.5-g dose; there was no further increase at higher doses.''
\noindent\hspace*{2em}\small\textit{Evidence:} ``In conclusion, the pharmacokinetics of oral L-carnitine display nonlinearity above a dose of 0.5 g 3 times a day.''
\noindent [272] BulkSupplements L-Carnitine L-Tartrate Powder 1 Kilogram - Target. \textit{Target}. \url{https://www.target.com/p/bulksupplements-l-carnitine-l-tartrate-powder-1-kilogram-2-2-lbs/-/A-1004332178} (accessed 2026-11-15).
\noindent\hspace*{2em}\small\textit{Evidence:} ``BulkSupplements L-Carnitine L-Tartrate Powder 1 Kilogram 0 out of 5 stars be the first! \$65.97''
\noindent [273] Nutricost L-Carnitine Tartrate Capsules. \textit{Nutricost}. \url{https://nutricost.com/products/nutricost-l-carnitine-tartrate-500mg-240-capsules} (accessed 2026-11-15).
\noindent\hspace*{2em}\small\textit{Evidence:} ``120/240 Capsules, Each With 500mg of L-Carnitine Tartrate''
\noindent\hspace*{2em}\small\textit{Evidence:} ``1000mg Per Serving \& 60/120 Servings Per Bottle''
\noindent [275] van Rijt WJ; Schielen PCJI; Özer Y; Bijsterveld K; van der Sluijs FH; Derks TGJ; Heiner-Fokkema MR (2020). Instability of Acylcarnitines in Stored Dried Blood Spots: The Impact on Retrospective Analysis of Biomarkers for Inborn Errors of Metabolism.. \textit{Int J Neonatal Screen} 6(4):83. \href{https://doi.org/10.3390/ijns6040083}{doi:10.3390/ijns6040083}. \url{https://pubmed.ncbi.nlm.nih.gov/33147805} (accessed 2026-11-15).
\noindent\hspace*{2em}\small\textit{Evidence:} ``Acylcarnitine species are subject to degradation upon long-term storage at room temperature, but limited data are available on the stability in original samples and the impact on acylcarnitine ratios.''
\noindent [276] Zhang Y; Jiang H; Hutson P (2012). Stability of acetyl-1-carnitine in 5\% dextrose using a high-performance liquid chromatography-mass spectrometry times 2 method.. \textit{Int J Pharm Compd} 16(2):170-3. \href{https://pubmed.ncbi.nlm.nih.gov/23050330/}{PMID:23050330}. \url{https://pubmed.ncbi.nlm.nih.gov/23050330} (accessed 2026-11-15).
\noindent\hspace*{2em}\small\textit{Evidence:} ``Acetyl-l-carnitine dissolved in aqueous solutions is stable at neutral to acidic pH, but unstable at pH > 9.''
\noindent\hspace*{2em}\small\textit{Evidence:} ``Acetyl-l-carnitine dissolved in water (pH 5.2) at concentrations of 1 and 10 mg/mL was found stable at room temperature or refrigerated for at least 33 days using the established stability-indicating method.''
